%% Question bank chapter file - generated by bank_tools.py build (do not edit while the bank is still being built:
%% it is regenerated from the per-test data in Question Bank\data). Every question is one bankq block:
%%   \begin{bankq}[id=..., source={...}, paper=.., qno=.., type=..., subtopic={...}, answer={...}, answer_src=key|solved,
%%                 working={...}, difficulty=1-10, diff_parts={K S I C F}, diff_note={...},
%%                 group=<passage id or empty>, also_in={duplicate copies in other tests}]
%% difficulty = min(10, K+S+I+C+F-1): K concepts 1-3, S steps 1-4, I insight 0-2, C computation 0-2, F format 0-1
%%   ...question LaTeX, options, figures (\bfig / \bfigi -> img/)...
%%   \end{bankq}
%% Passages shared by several questions are bankpassage blocks; their questions carry group=<passage id>.
\documentclass[10pt]{article}
\usepackage[a4paper,margin=1.8cm]{geometry}
\usepackage[T1]{fontenc}
\usepackage{amsmath,amssymb,mathtools}
\usepackage{graphicx}
\usepackage[export]{adjustbox}
\usepackage{xcolor,needspace,keyval}
\usepackage{enumitem}
\usepackage{booktabs,array,tabularx}
\usepackage[hidelinks]{hyperref}
\setlength{\parindent}{0pt}
\setlength{\parskip}{3pt}
\makeatletter
\define@key{bq}{id}{\def\bq@id{#1}}\define@key{bq}{source}{\def\bq@src{#1}}\define@key{bq}{paper}{\def\bq@paper{#1}}
\define@key{bq}{qno}{\def\bq@qno{#1}}\define@key{bq}{type}{\def\bq@type{#1}}\define@key{bq}{subtopic}{\def\bq@sub{#1}}
\define@key{bq}{answer}{\def\bq@ans{#1}}\define@key{bq}{answer_src}{\def\bq@asrc{#1}}\define@key{bq}{working}{\def\bq@work{#1}}
\define@key{bq}{group}{\def\bq@group{#1}}\define@key{bq}{also_in}{\def\bq@also{#1}}
\define@key{bq}{difficulty}{\def\bq@diff{#1}}\define@key{bq}{diff_parts}{\def\bq@dparts{#1}}\define@key{bq}{diff_note}{\def\bq@dnote{#1}}
\newenvironment{bankq}[1][]{\def\bq@work{}\def\bq@also{}\def\bq@group{}\def\bq@diff{?}\def\bq@dparts{}\def\bq@dnote{}\setkeys{bq}{#1}%
  \par\needspace{4\baselineskip}\medskip\noindent{\small\textbf{\bq@id}\hfill\color{gray}[\bq@src, Paper \bq@paper, Q\bq@qno]
  \bq@type\ $\cdot$ \bq@sub\ $\cdot$ \textbf{D\bq@diff}}\par\nopagebreak\relax}
  {\par{\small\color{teal}Answer: \textbf{\bq@ans}\ifx\bq@work\@empty\else\ (solved: \bq@work)\fi
  \ifx\bq@also\@empty\else\par\color{gray}Also in: \bq@also\fi}\par}
\define@key{bp}{id}{\def\bp@id{#1}}\define@key{bp}{source}{\def\bp@src{#1}}\define@key{bp}{questions}{\def\bp@qs{#1}}
\newenvironment{bankpassage}[1][]{\setkeys{bp}{#1}\par\needspace{4\baselineskip}\bigskip\noindent
  {\small\textbf{Passage \bp@id}\hfill\color{gray}[\bp@src; questions \bp@qs]}\par\nopagebreak\relax\itshape}{\par}
\makeatother
\newcommand{\bfig}[2]{\par\smallskip{\centering\includegraphics[width=#1pt,max width=0.62\linewidth,max height=0.32\textheight]{img/#2}\par}\smallskip}
\newcommand{\bfigi}[2]{\includegraphics[width=#1pt,max width=\linewidth,max height=0.32\textheight,valign=c]{img/#2}}
\newcommand{\bmissing}[1]{\textit{[figure missing in the source conversion: \detokenize{#1}]}}
\newenvironment{qoptions}{\begin{itemize}[leftmargin=1.8em,itemsep=1pt,topsep=2pt,label={}]}{\end{itemize}}
\newcommand{\opt}[2]{\item[(#1)] #2}
\providecommand{\fwstart}{}\providecommand{\fwend}{}\providecommand{\qfullsplit}{}
\newcolumntype{L}[1]{>{\raggedright\arraybackslash}p{#1}}
\newcolumntype{C}[1]{>{\centering\arraybackslash}p{#1}}
\newcolumntype{R}[1]{>{\raggedleft\arraybackslash}p{#1}}
\providecommand{\sectionhdr}[1]{}\providecommand{\subsectionhdr}[1]{}


\begin{document}
{\Large\bfseries Trig Functions \& Equations --- Question Bank}\par\smallskip
{\small Mathematics chapter M12. Questions from all converted test papers, grouped by source series. Answers: printed key where it matches the paper, otherwise solved (answer\_src=solved, with working). Exact repeats across papers appear once, with also\_in listing the other copies.}\par\medskip\hrule

\section*{Score-I papers}
\begin{bankq}[id=SI04-P2-MATH-S1-Q06, source={Score-I - Test 04 (14 Jan 2022)}, paper=2, qno=6, type=single, subtopic={General solutions}, answer={C}, answer_src=key, difficulty=6, diff_parts={K2 S3 I1 C1 F0}, diff_note={concepts: number, solutions of exponential trigonometric equation | steps: let t = 2^{sin^2 x}; 4 times 2^{1 - sin^2x}; quadratic in t; solve; count x in the interval}]
Number of values of `$x$' in $(-2\pi,2\pi)$ satisfying the equation $2^{\sin^2x}+4\times2^{\cos^2x}=6$ is -
\begin{qoptions}
\opt{A}{8}
\opt{B}{6}
\opt{C}{4}
\opt{D}{2}
\end{qoptions}
\end{bankq}
\begin{bankq}[id=SI06-P1-MATH-S1-Q03, source={Score-I - Test 06 (28 Jan 2022)}, paper=1, qno=3, type=single, subtopic={General solutions}, answer={C}, answer_src=key, difficulty=3, diff_parts={K1 S1 I1 C1 F0}, diff_note={concepts: which equation has real solutions | steps: bounds on trigonometric terms or algebraic discriminants; check each equation}]
Which of the following equation has real solution(s) ?
\begin{qoptions}
\opt{A}{$\cos(\sin x)=\dfrac12$}
\opt{B}{$(x^2+1)^{10}=2x-x^2-2$}
\opt{C}{$(3+\cos x)^2=4-2\sin^2x$}
\opt{D}{$\sin x+\cos x=\dfrac{\pi}{2}$}
\end{qoptions}
\end{bankq}
\begin{bankq}[id=SI06-P1-MATH-S3-Q01, source={Score-I - Test 06 (28 Jan 2022)}, paper=1, qno=1, type=numerical, subtopic={Graphs}, answer={1.00}, answer_src=key, difficulty=5, diff_parts={K2 S1 I1 C1 F1}, diff_note={concepts: number, solutions of a nested trigonometric equation | steps: compare sin(cos(sin x)) and cos(sin(cos x)) over the interval using monotonic behaviour; count}]
The number of solutions that the equation $\sin(\cos(\sin x))=\cos(\sin(\cos x))$ has in $\left[0,\dfrac{\pi}{2}\right]$ is :
\end{bankq}
\begin{bankq}[id=SI10-P1-MATH-S2-Q05, source={Score-I - Test 10 (24 Feb 2022)}, paper=1, qno=5, type=numerical, subtopic={General solutions}, answer={4}, answer_src=key, difficulty=5, diff_parts={K1 S2 I1 C1 F1}, diff_note={concepts: trigonometric equation as a polynomial in sin x | steps: take s = sin x and factor 4s^4+4s^3+3s^2-s-1 using s=1/2 and s=-1/2; remaining quadratic has no real root; count solutions of sin x = 1/2 and -1/2 in the interval | I1: finding rational roots of the quartic}]
The number of distinct real solutions of equation $4\sin^4x+4\sin^3x+3\sin^2x=1+\sin x$, in the interval $[0,2\pi]$, is
\end{bankq}
\begin{bankq}[id=SI11-P1-MATH-S3-Q03, source={Score-I - Test 11 (27 Feb 2022)}, paper=1, qno=3, type=numerical, subtopic={Trig inequalities}, answer={1.00}, answer_src=key, difficulty=6, diff_parts={K2 S2 I1 C1 F1}, diff_note={concepts: length, the interval where a trigonometric function is nonnegative | steps: use product to sum on f; reduce to single sine expression; solve inequality; difference of endpoints}]
Let $f:[0,2]\to\mathbb R$ be the function defined by $f(x)=(3-\sin(2\pi x))\sin\left(\pi x-\dfrac\pi4\right)-\sin\left(3\pi x+\dfrac\pi4\right)$. If $\alpha,\beta\in[0,2]$ are such that $\{x\in[0,2]:f(x)\ge0\}=[\alpha,\beta]$, then the value of $\beta-\alpha$ is \underline{\hspace{1.2cm}}
\end{bankq}
\begin{bankq}[id=SI18-P1-MATH-S3-Q06, source={Score-I - Test 18 (24 Apr 2022)}, paper=1, qno=6, type=numerical, subtopic={General solutions}, answer={12.57}, answer_src=solved, working={By AM-GM the left side is at least 4 while the right side is at most 4; equality needs tan^2 x = tan^2 y = 1 and sin^2(x+y) = 1, so x takes the values pi/4, 3pi/4, 5pi/4, 7pi/4 and their sum is 4 pi = 12.57. The printed key letters for this row do not fit the question.}, difficulty=7, diff_parts={K2 S3 I1 C1 F1}, diff_note={concepts: sum, x values satisfying a trigonometric equation with tangents | steps: let a = tan^2 x and b = tan^2 y; AM-GM gives left side at least 4; right side at most 4; equality forces a = b = 1 and sin^2(x+y) = 1; list x values in interval; add}]
If $(x,y)$ satisfy the equation $\tan^4x+\tan^4y+2\cot^2x\cot^2y=3+\sin^2(x+y)$ sum of all the possible values of $x$ where $x,y\in(0,2\pi)$
\end{bankq}
\begin{bankq}[id=SI19-P1-MATH-S2-Q03, source={Score-I - Test 19 (20 Nov 2022)}, paper=1, qno=3, type=numerical, subtopic={General solutions}, answer={0.00}, answer_src=key, difficulty=7, diff_parts={K2 S3 I1 C1 F1}, diff_note={concepts: number, real solutions of a trigonometric equation with maximum bounds | steps: let t = sin theta + cos theta; sin 2 theta = t^2 - 1; equation becomes polynomial in t; range of t restricts; count theta in [0, 4 pi]}]
Number of real solution(s) of the equation $\dfrac{\sin^22\theta}{2}+6\sin\theta\cdot\cos\theta(\sin\theta+\cos\theta)=10$ in $[0,4\pi]$ is
\end{bankq}
\begin{bankq}[id=SI19-P1-MATH-S2-Q04, source={Score-I - Test 19 (20 Nov 2022)}, paper=1, qno=4, type=numerical, subtopic={General solutions}, answer={6.00}, answer_src=key, difficulty=6, diff_parts={K2 S2 I1 C1 F1}, diff_note={concepts: number of angle values from a logarithmic condition on p, q | steps: from (log_p q + log_q p)^2 = 2029 obtain (log_q p - log_p q)^2 = 2025 so difference 45; then cosec theta = 45 or related; count solutions in (0, 6 pi)}]
If $\left(\log_pq+\log_qp\right)^2=2029$, $p>q>1$, then the number of values of $\theta\in(0,6\pi)$ satisfying the equation $\dfrac{1}{\log_{pq}q}-\dfrac{1}{\log_{pq}p}=\operatorname{cosec}\theta$ is
\end{bankq}
\begin{bankq}[id=SI19-P1-MATH-S3-Q02, source={Score-I - Test 19 (20 Nov 2022)}, paper=1, qno=2, type=numerical, subtopic={Trig inequalities}, answer={5}, answer_src=key, difficulty=6, diff_parts={K2 S2 I1 C1 F1}, diff_note={concepts: cardinality, the set built from three trigonometric conditions | steps: solve each of the three conditions in its interval; union of A and B; intersect with inequality set C; count}]
Consider the following sets\\
$A=\{\theta:\sec\theta+\operatorname{cosec}\theta=\sin\theta+\cos\theta,\ \theta\in[0,2\pi]\}$,\\
$B=\left\{\theta:\cos\theta=\sqrt{1-\sin\theta},\ \theta\in[0,3\pi]\right\}$,\\
$C=\Bigl\{\theta:2^{\sin^2\theta}>\dfrac{\cos\theta}{2\cos\theta+1},$\\
\hspace*{1em}$\theta\in\mathbb R\Bigr\}$. The number of elements in the set $(A\cup B)\cap C$ is
\end{bankq}
\begin{bankq}[id=SI19-P1-MATH-S3-Q03, source={Score-I - Test 19 (20 Nov 2022)}, paper=1, qno=3, type=numerical, subtopic={General solutions}, answer={3}, answer_src=key, difficulty=6, diff_parts={K2 S2 I1 C1 F1}, diff_note={concepts: number, solutions of a logarithmic trigonometric equation | steps: combine logs to cos^3 5x + cos^3 7x = 8 cos 6x cos x?; identity a^3 + b^3 with sum 2 cos 6x cos x; solve; count in the interval}]
Number of real solution(s) of equation $\dfrac13\log_2(\cos^35x+\cos^37x)=1+\log_2\cos6x+\log_2\cos x$ in $\left[0,\dfrac\pi2\right]$ is
\end{bankq}
\begin{bankq}[id=SI31-P1-MATH-S2-Q04, source={Score-I - Test 31 (29 Dec 2022)}, paper=1, qno=4, type=numerical, subtopic={General solutions}, answer={4.00}, answer_src=key, difficulty=6, diff_parts={K1 S2 I2 C1 F1}, diff_note={concepts: number of solutions of a trigonometric equation with tan sin cot cos in the interval | steps: rewrite using sin and cos; factor common terms into (sin - cos) forms; substitute t = sin + cos or solve factored equation; count roots in 0 to 2 pi | I: grouping terms reveals a factor}]
The number of solutions of the equation $2(\tan x-\sin x)+3(\cot x-\cos x)+5=0$ in $x\in(0,2\pi)$ is
\end{bankq}
\begin{bankq}[id=SI31-P2-MATH-S1-Q01, source={Score-I - Test 31 (29 Dec 2022)}, paper=2, qno=1, type=single, subtopic={General solutions}, answer={C}, answer_src=key, difficulty=5, diff_parts={K1 S2 I2 C1 F0}, diff_note={concepts: number of elements in a set of angles satisfying a tan equation with double angle | steps: substitute t = tan theta and tan 2 theta formula; reduce to equation in t; solve and list angles in the set S excluding forbidden points; count | I: remove excluded angles from the solution set}]
Let $S=\left[-\pi,\dfrac{\pi}{2}\right)-\left\{-\dfrac{\pi}{2},-\dfrac{\pi}{4},-\dfrac{3\pi}{4},\dfrac{\pi}{4}\right\}$.

Then the number of elements in the set

$A=\left\{\theta\in S:\tan\theta\left(1+\sqrt3\tan2\theta\right)=\sqrt3-\tan2\theta\right\}$ is
\begin{qoptions}
\opt{A}{3}
\opt{B}{4}
\opt{C}{5}
\opt{D}{7}
\end{qoptions}
\end{bankq}
\begin{bankq}[id=SI31-P2-MATH-S2-Q01, source={Score-I - Test 31 (29 Dec 2022)}, paper=2, qno=1, type=single, subtopic={General solutions}, answer={5}, answer_src=key, difficulty=4, diff_parts={K1 S2 I1 C1 F0}, diff_note={concepts: number solutions of a mixed trigonometric equation | steps: use product to sum for cos x sin 2x; combine terms to factor (1 + sin x)(cos 2x + cos x) = 0; solve each factor in the interval; count 5}]
$2(\cos x+\cos2x)+(1+2\cos x)\sin2x=2\sin x$

Total no. of solutions in $[-\pi,\pi]$ is
\end{bankq}

\section*{Score-II papers}
\begin{bankq}[id=SII02-P1-MATH-S3-Q03, source={Score-II - Test 02 (28 Mar 2021)}, paper=1, qno=3, type=numerical, subtopic={General solutions}, answer={4}, answer_src=key, difficulty=5, diff_parts={K1 S2 I1 C1 F1}, diff_note={concepts: trigonometric equation with squared reciprocal functions | steps: express everything in sin and cos; combine to a single trigonometric equation; solve; count in the interval}]
Find the number of solutions of $\theta$ in $(0,2\pi)$ that satisfy the equation
\[\dfrac{1}{\sin^2\theta}-\dfrac{1}{\cos^2\theta}-\dfrac{1}{\tan^2\theta}-\dfrac{1}{\cot^2\theta}-\dfrac{1}{\sec^2\theta}-\dfrac{1}{\operatorname{cosec}^2\theta}=-3\]
\end{bankq}
\begin{bankq}[id=SII03-P1-MATH-S1-Q08, source={Score-II - Test 03 (05 May 2021)}, paper=1, qno=8, type=single, subtopic={General solutions}, answer={D}, answer_src=key, difficulty=3, diff_parts={K1 S2 I1 C0 F0}, diff_note={concepts: ordered pairs satisfying x^4 - 2x^2 cos t + 1 = 0 | steps: discriminant condition cos^2 t >= 1; cos t = plus minus 1; solutions of x^2; count pairs}]
If $\theta\in[0,2\pi]$ then number of orderered pairs (x, $\theta$) which satisfies $x^4-2x^2\cos\theta+1=0$ is/are
\begin{qoptions}
\opt{A}{1}
\opt{B}{2}
\opt{C}{3}
\opt{D}{4}
\end{qoptions}
\end{bankq}
\begin{bankq}[id=SII05-P1-MATH-S1-Q12, source={Score-II - Test 05 (22 Mar 2022)}, paper=1, qno=12, type=single, subtopic={Trig inequalities}, answer={D}, answer_src=key, difficulty=4, diff_parts={K1 S2 I1 C1 F0}, diff_note={concepts: interval satisfying a trigonometric and a quadratic inequality | steps: factor 2 sin^2 x + 3 sin x - 2; sin x greater than 1/2; second inequality gives -1 < x < 2; intersect}]
Let $2\sin^2x+3\sin x-2>0$ and $x^2-x-2<0$ (x is in radians), then x lies in the interval $-$
\begin{qoptions}
\opt{A}{$\left(\dfrac\pi6,\dfrac{5\pi}6\right)$}
\opt{B}{$\left(-1,\dfrac{5\pi}6\right)$}
\opt{C}{$(-1,2)$}
\opt{D}{$\left(\dfrac\pi6,2\right)$}
\end{qoptions}
\end{bankq}
\begin{bankq}[id=SII07-P1-MATH-S1-Q09, source={Score-II - Test 07 (10 Jul 2022)}, paper=1, qno=9, type=single, subtopic={General solutions}, answer={D}, answer_src=key, difficulty=3, diff_parts={K1 S2 I1 C0 F0}, diff_note={concepts: product of three linear factors in sin t equal to 6 | steps: product increasing in sin t; equals 6 at sin t = -1; count solutions in the interval}]
The number of solutions of the equation $\sqrt2\cos3x-\cos2x=\sqrt3\sin2x+\sqrt2\sin3x$ in $\left[\dfrac{-4\pi}{3},\dfrac{2\pi}{3}\right]$ is
\begin{qoptions}
\opt{A}{3}
\opt{B}{4}
\opt{C}{5}
\opt{D}{6}
\end{qoptions}
\end{bankq}
\begin{bankq}[id=SII09-P1-MATH-S3-Q02, source={Score-II - Test 09 (12 Feb 2023)}, paper=1, qno=2, type=numerical, subtopic={General solutions}, answer={Bonus}, answer_src=solved, working={Printed key for this numeric answer is corrupted (a multiple-choice letter pattern) and the question was not solved; banked as Bonus.}, difficulty=4, diff_parts={K1 S1 I2 C0 F1}, diff_note={concepts: trigonometric equation by Cauchy-Schwarz | steps: expression is a combination with maximum value sqrt of sum of squares which equals 7 only for equality; equality conditions give tan x and tan y | I: the left side is bounded by the Cauchy-Schwarz amplitude 7 so equality fixes both angles}]
If $x,y\in(0,2\pi)$, satisfying: $2\sin x\cdot\sin y+3\cos y+6\cos x\cdot\sin y=7$ then the value of $(\tan^2x+\tan^2y)$ is
\end{bankq}
\begin{bankq}[id=SII12-P1-MATH-S4-Q01, source={Score-II - Test 12 (04 Feb 2024)}, paper=1, qno=1, type=numerical, subtopic={General solutions}, answer={3}, answer_src=key, difficulty=5, diff_parts={K1 S2 I1 C1 F1}, diff_note={concepts: solutions of an equation with inverse tangent of tangent | steps: principal value reduces the right side to a piecewise linear function; the left side is a cosine form; solve in each interval; count solutions}]
Let $2\tan^{-1}x\in(-\pi,\pi)\ \forall x\in R$. The number of real solutions of\\
\resizebox{\linewidth}{!}{$\sqrt{1+\cos\dfrac{2x}{3}}=\sqrt2\tan^{-1}\tan\dfrac x3;\ x\in\left(-\dfrac{9\pi}{2},\dfrac{9\pi}{2}\right);\ x\ne\dfrac{(6n+3)\pi}{2},\ n\in Z$}\\
is equal to
\end{bankq}
\begin{bankq}[id=SII17-P1-MATH-S1-Q18, source={Score-II - Test 17 (13 Feb 2024)}, paper=1, qno=18, type=single, subtopic={General solutions}, answer={A}, answer_src=solved, working={u=9\^{}t: u+9/u=10: tan\^{}2x=0 or 1: x=0,+-pi/4; beta=2tan\^{}2(pi/12)=14-8sqrt3; (beta-14)\^{}2/6=192/6=32.}, difficulty=5, diff_parts={K1 S3 I1 C1 F0}, diff_note={concepts: sum over solutions of an exponential equation in tan x | steps: set u = 9^(tan^2 x): 9/u + u = 10; u = 1 or 9; tan^2 x = 0 or 1; sum of tan^2(x/3) over the solutions; evaluate 16(beta - )}]
Let $S=\left\{x\in\left(-\dfrac\pi2,\dfrac\pi2\right):9^{1-\tan^2x}+9^{\tan^2x}=10\right\}$ and $\beta=\displaystyle\sum_{x\in S}\tan^2\left(\dfrac x3\right)$, then $\dfrac16(\beta-14)^2$ is equal to
\begin{qoptions}
\opt{A}{32}
\opt{B}{8}
\opt{C}{64}
\opt{D}{16}
\end{qoptions}
\end{bankq}
\begin{bankq}[id=SII19-P1-MATH-S2-Q07, source={Score-II - Test 19 (20 Feb 2024)}, paper=1, qno=7, type=numerical, subtopic={General solutions}, answer={4}, answer_src=key, difficulty=6, diff_parts={K1 S3 I1 C1 F1}, diff_note={concepts: trigonometric identity with powers of sine | steps: sin^2 a + sin^4 a = 1 gives sin^2 a = (sqrt5 - 1)/2; cos^2 a = sin^4 a; compute cos^8 - sin^8; match with b - 9 sqrt a; b a}]
If $\sin^2\alpha+\sin^4\alpha=1$ \& $\cos^8\alpha-\sin^8\alpha=\left(b-9\sqrt a\right)$, then $\dfrac ba$ is equal to
\end{bankq}
\begin{bankq}[id=SII20-P1-MATH-S1-Q17, source={Score-II - Test 20 (23 Feb 2024)}, paper=1, qno=17, type=single, subtopic={General solutions}, answer={C}, answer_src=key, difficulty=3, diff_parts={K1 S2 I1 C0 F0}, diff_note={concepts: product of three linear factors in sin t equal to 6 | steps: product increasing in sin t; equals 6 at sin t = -1; count solutions in the interval}]
Number of solutions of $5\cos^2\theta-3\sin^2\theta+6\sin\theta\cos\theta=7$ in the interval $[0,2\pi]$ is :-
\begin{qoptions}
\opt{A}{2}
\opt{B}{4}
\opt{C}{0}
\opt{D}{None of these}
\end{qoptions}
\end{bankq}
\begin{bankq}[id=SII20-P1-MATH-S2-Q09, source={Score-II - Test 20 (23 Feb 2024)}, paper=1, qno=9, type=numerical, subtopic={General solutions}, answer={Bonus}, answer_src=solved, working={Printed key omits or corrupts this answer (letters shown for numeric-type items) and the question was not solved; banked as Bonus.}, difficulty=6, diff_parts={K1 S2 I2 C1 F1}, diff_note={concepts: number of x in an interval with an equation equal 2 | steps: each term bounded so equality needs sin x = 1 and cos 2x = -1; check consistency; count over 10 pi | I: extremal values force a specific point}]
The number of values of $x$ in $[0,10\pi]$ for which $(\sin x)^{15}-(\cos2x)^{19}=2$ is
\end{bankq}
\begin{bankq}[id=SII20-P2-MATH-S4-Q06, source={Score-II - Test 20 (23 Feb 2024)}, paper=2, qno=6, type=numerical, subtopic={General solutions}, answer={4}, answer_src=key, difficulty=4, diff_parts={K1 S1 I1 C1 F1}, diff_note={concepts: number of solutions of a trigonometric equation on (0  4 pi) | steps: reduce to simple trig equation using identities; count in each period}]
The number of solution satisfying the equation $\begin{vmatrix}1+\sin^3\theta&1\\-\cos^3\theta&1\end{vmatrix}=\begin{vmatrix}\cos\theta&4\cos\theta\\0&3\sin\theta\end{vmatrix}$ in $\theta\in(0,4\pi)$ is equal to
\end{bankq}
\begin{bankq}[id=SII25-P1-MATH-S1-Q18, source={Score-II - Test 25 (20 Mar 2024)}, paper=1, qno=18, type=single, subtopic={General solutions}, answer={D}, answer_src=key, difficulty=5, diff_parts={K1 S2 I2 C1 F0}, diff_note={concepts: sum of solutions of a trigonometric equation with inverse functions | steps: simplify arctan cot of an angle using reduction; sin theta equals sin of the value; solutions in range; sum | I: reduce the inverse function expression to the principal value}]
If $S$ represent the sum of the solutions of the equation $\sin\theta=\sin\left(\tan^{-1}\cot\left(\dfrac{75\pi}{8}\right)\right)$ in $[0,2\pi]$, then $\cos S$ is -
\begin{qoptions}
\opt{A}{$\dfrac12\sqrt{2+\sqrt2}$}
\opt{B}{$\dfrac12\sqrt{2-\sqrt2}$}
\opt{C}{1}
\opt{D}{$-1$}
\end{qoptions}
\end{bankq}
\begin{bankq}[id=SII27-P1-MATH-S2-Q01, source={Score-II - Test 27 (23 Mar 2024)}, paper=1, qno=1, type=numerical, subtopic={General solutions}, answer={0}, answer_src=key, difficulty=5, diff_parts={K1 S2 I1 C1 F1}, diff_note={concepts: number of beta satisfying a trigonometric identity | steps: simplify the radicand 2 cot a + 1/sin^2 a = (cot a + 1)^2; so sqrt tan beta =; beta determined; count}]
Find number of possible values of $\beta$, if $\sqrt{2\cot\alpha+\dfrac1{\sin^2\alpha}}=\dfrac{1+\cot\alpha}{\sqrt{\tan\beta}}$ for $\alpha\in\left(\dfrac{3\pi}4,\pi\right)$
\end{bankq}
\begin{bankq}[id=SII30-P1-MATH-S1-Q11, source={Score-II - Test 30 (27 Mar 2024)}, paper=1, qno=11, type=single, subtopic={General solutions}, answer={D}, answer_src=key, difficulty=4, diff_parts={K1 S2 I1 C1 F0}, diff_note={concepts: trigonometric equation as a polynomial in sin x | steps: take s = sin x and factor 4s^4+4s^3+3s^2-s-1 using s=1/2 and s=-1/2; remaining quadratic has no real root; count solutions of sin x = 1/2 and -1/2 in the interval | I1: finding rational roots of the quartic}]
The number of solutions of the equation $4\sin^2x-4\cos^3x+9-4\cos x=0$; $x\in[-2\pi,2\pi]$ is:
\begin{qoptions}
\opt{A}{1}
\opt{B}{3}
\opt{C}{2}
\opt{D}{0}
\end{qoptions}
\end{bankq}
\begin{bankq}[id=SII36-P1-MATH-S1-Q02, source={Score-II - Test 36 (07 Feb 2025)}, paper=1, qno=2, type=single, subtopic={Trig inequalities}, answer={D}, answer_src=key, difficulty=3, diff_parts={K1 S2 I1 C0 F0}, diff_note={concepts: amplitude bound of a trigonometric inequality for all x | steps: sin x - cos(alpha) cos x has minimum -sqrt(1+cos^2 alpha); needs sin y at least sqrt(1+cos^2 alpha) which forces sin y=1 and cos alpha=0; sum 1}]
If $\sin x+\sin y\geq\cos\alpha\cdot\cos x,\ \forall x\in\mathbb R$, then $\sin y+\cos\alpha$ is equal to

\begin{qoptions}
\opt{A}{$-2$}
\opt{B}{$-\dfrac32$}
\opt{C}{0}
\opt{D}{1}
\end{qoptions}
\end{bankq}
\begin{bankq}[id=SII39-P2-MATH-S1-Q01, source={Score-II - Test 39 (02 Mar 2025)}, paper=2, qno=1, type=single, subtopic={General solutions}, answer={A}, answer_src=key, difficulty=5, diff_parts={K1 S3 I1 C1 F0}, diff_note={concepts: sum of solutions of a palindromic quartic in cotangent | steps: divide by cot^2 and put t = cot + 1/cot; quadratic in t; solve; find cot values; angles in (0,2pi) excluding pi; sum}]
For $\theta\in(0,2\pi)-\{\pi\}$, sum of all values of $`\theta'$ satisfying equation $4\cot^4\theta+9\cot^3\theta+8\cot^2\theta+9\cot\theta+4=0$; is equal to :
\begin{qoptions}
\opt{A}{$5\pi$}
\opt{B}{$6\pi$}
\opt{C}{$7\pi$}
\opt{D}{None}
\end{qoptions}
\end{bankq}
\begin{bankq}[id=SII40-P1-MATH-S4-Q01, source={Score-II - Test 40 (16 Mar 2025)}, paper=1, qno=1, type=numerical, subtopic={General solutions}, answer={9}, answer_src=key, difficulty=5, diff_parts={K1 S2 I1 C1 F1}, diff_note={concepts: solutions of an equation with inverse tangent of tangent | steps: principal value reduces the right side to a piecewise linear function; the left side is a cosine form; solve in each interval; count solutions}]
Let $\tan^{-1}(x)\in\left(-\dfrac\pi2,\dfrac\pi2\right)$, for $x\in\mathbb R$. Then the number of real solutions of the equation\\
$\sqrt{1+\cos(2x)}=\sqrt2\tan^{-1}\left(\tan\dfrac x6\right)$ in the set $(-9\pi,-3\pi)\cup(-3\pi,3\pi)\cup(3\pi,9\pi)$ is equal to
\end{bankq}
\begin{bankq}[id=SII41-P1-MATH-S2-Q03, source={Score-II - Test 41 (19 Mar 2025)}, paper=1, qno=3, type=numerical, subtopic={General solutions}, answer={4}, answer_src=key, difficulty=5, diff_parts={K1 S2 I1 C1 F1}, diff_note={concepts: trigonometric equation as a polynomial in sin x | steps: take s = sin x and factor 4s^4+4s^3+3s^2-s-1 using s=1/2 and s=-1/2; remaining quadratic has no real root; count solutions of sin x = 1/2 and -1/2 in the interval | I1: finding rational roots of the quartic}]
Number of roots of the equation $\cos^2x+\dfrac{\left(\sqrt3+1\right)}{2}\sin x-\dfrac{\sqrt3}4-1=0$ which lie in the interval $[-\pi,\pi]$ is
\end{bankq}
\begin{bankq}[id=SII42-P1-MATH-S3-Q01, source={Score-II - Test 42 (19 Mar 2025)}, paper=1, qno=1, type=numerical, subtopic={General solutions}, answer={2}, answer_src=key, difficulty=6, diff_parts={K2 S2 I1 C1 F1}, diff_note={concepts: critical value of a trigonometric equation, integer roots of a quadratic | steps: c = 2t/(t^2-1) with t = sin+cos in [-sqrt2,sqrt2]; find the range of c^2 where the number of solutions changes to get k; then count integer solutions of x^2+6x+k=0}]
The equation $\sec\theta+\text{cosec}\theta=c$ has two solutions between 0 and $2\pi$ if $c^2<k$ and four solutions if $c^2>k$. Find number of integral solution of the equation $x^2+k+6x=0$
\end{bankq}
\begin{bankq}[id=SII43-P1-MATH-S1-Q05, source={Score-II - Test 43 (23 Mar 2025)}, paper=1, qno=5, type=single, subtopic={General solutions}, answer={A}, answer_src=key, difficulty=4, diff_parts={K1 S2 I1 C1 F0}, diff_note={concepts: least n for seven solutions of trig equation | steps: convert to quadratic in sec; accept the valid root; list solutions in each quadrant; find n}]
If $2\tan^2\theta-5\sec\theta=1$ has exactly 7 solutions in the interval $\left[0,\dfrac{n\pi}2\right]$, then for the least value of $n\in N$ value of $\displaystyle\sum_{k=1}^n\dfrac k{2^k}$ is equal to
\begin{qoptions}
\opt{A}{$\dfrac1{2^{13}}\left(2^{14}-15\right)$}
\opt{B}{$1-\dfrac{15}{2^{13}}$}
\opt{C}{$\dfrac1{2^{15}}\left(2^{14}-14\right)$}
\opt{D}{$\dfrac1{2^{14}}\left(2^{15}-15\right)$}
\end{qoptions}
\end{bankq}
\begin{bankq}[id=SII44-P1-MATH-S1-Q03, source={Score-II - Test 44 (23 Mar 2025)}, paper=1, qno=3, type=single, subtopic={Graphs}, answer={Bonus}, answer_src=solved, working={Printed key covers only one subject (chemistry) and this question was not solved; banked as Bonus.}, difficulty=6, diff_parts={K2 S3 I1 C1 F0}, diff_note={concepts: inverse trigonometric identities, common tangent to an ellipse and a parabola | steps: for positive x the sum is pi/2+pi/2 giving k=1 and for negative x it gives k=0; ellipse and parabola from k1 and k2; tangent y = m x + c to the parabola and condition for the ellipse; solve for m c; compare}]
If $\sin^{-1}x+\cos^{-1}x+\tan^{-1}\dfrac1x+\tan^{-1}x=\pi k$. The value of k for different values of x are k$_1$ and k$_2$, satisfying the equations $\dfrac{x^2}{4k_1}+\dfrac{y^2}5=1$ and $y^2=12(x-k_2)$, then the equation of their common tangent is

\begin{qoptions}
\opt{A}{$y=2x-5$}
\opt{B}{$3y=x+3$}
\opt{C}{$y=x+3$}
\opt{D}{$y=3x+1$}
\end{qoptions}
\end{bankq}
\begin{bankq}[id=SII45-P1-MATH-S1-Q16, source={Score-II - Test 45 (26 Mar 2025)}, paper=1, qno=16, type=single, subtopic={General solutions}, answer={C}, answer_src=key, difficulty=6, diff_parts={K1 S3 I2 C1 F0}, diff_note={concepts: find x and y from a quadratic in log y | steps: treat as quadratic in log y; discriminant by AM GM forces equality; x = 0; y value; log(x + y) | I: AM GM makes the discriminant zero only at a single point}]
If $\log^2y+\left(2^{1+x}+2^{1-x}\right)(\log y)+2^{2x+1}+2^{1-2x}=0$; where $x,y\in R$, then the value of $\log(x+y)$ is :- (where base of log is e)
\begin{qoptions}
\opt{A}{$-1$}
\opt{B}{1}
\opt{C}{$-2$}
\opt{D}{2}
\end{qoptions}
\end{bankq}

\section*{Score Advanced papers}
\begin{bankq}[id=SA01-P1-MATH-S1-Q06, source={Score Advanced - Test 01 (22 Sep 2020)}, paper=1, qno=6, type=single, subtopic={General solutions}, answer={B}, answer_src=key, difficulty=4, diff_parts={K1 S2 I1 C1 F0}, diff_note={concepts: trigonometric equation sum of roots | steps: convert to sin and cos; simplify to a quadratic in sin 2 theta; solve; collect roots}]
Sum of the solutions of the equation $\sin^2\theta\tan\theta+\cos^2\theta\cot\theta-\sin2\theta=1+\tan\theta+\cot\theta$, where $\theta\in(0,\pi)-\left\{\dfrac\pi2\right\}$, is -
\begin{qoptions}
\opt{A}{$3\pi$}
\opt{B}{$\dfrac{3\pi}2$}
\opt{C}{$\dfrac{5\pi}2$}
\opt{D}{$7\pi$}
\end{qoptions}
\end{bankq}
\begin{bankq}[id=SA08-P1-MATH-S3-Q02, source={Score Advanced - Test 08 (22 Apr 2023)}, paper=1, qno=2, type=single, subtopic={Graphs}, answer={8}, answer_src=key, difficulty=5, diff_parts={K1 S2 I2 C1 F0}, diff_note={concepts: equation with fractional part and inverse sine of sine | steps: piecewise form of sin^-1 sin x; fractional part; solve in each interval; count | I: periodic zigzag}]
Number of solution of the equation $\left\{\sin^{-1}(\sin\text{x})\right\}$ = $\dfrac{\text{x}}{10}$, in [0, 2$\pi$) is, where \{.\} denotes fractional part function :
\end{bankq}
\begin{bankq}[id=SA08-P2-MATH-S2-Q05, source={Score Advanced - Test 08 (22 Apr 2023)}, paper=2, qno=5, type=numerical, subtopic={General solutions}, answer={0.83}, answer_src=key, difficulty=6, diff_parts={K1 S2 I2 C1 F1}, diff_note={concepts: trigonometric equation with a quadratic bound | steps: range of 12 sin x + 5 cos x; quadratic in y minimum; equality case; solve | I: both sides at extreme values}]
If x and y are the solutions of the equation 12sin x + 5cos x = 2y$^2$ $-$ 8y + 21, the least positive value of $\dfrac{\text{x}}{\text{y}}$ is $\dfrac{1}{\text{a}}\cot^{-1}\left(\dfrac{\text{b}}{\text{c}}\right)$ then the value of $\left(\dfrac{\text{ab}}{\text{c}}\right)$ is :
\end{bankq}
\begin{bankq}[id=SA09-P1-MATH-S1-Q03, source={Score Advanced - Test 09 (29 Apr 2023)}, paper=1, qno=3, type=single, subtopic={Trig inequalities}, answer={C}, answer_src=key, difficulty=5, diff_parts={K1 S2 I2 C1 F0}, diff_note={concepts: number of ordered triplets satisfying a sum of fourth powers equals a product | steps: AM-GM on the three fourth powers with weights; equality occurs when sin values match specific numbers; equation forces sin x = sin y magnitudes; count angles in [0,2pi] for each sign pattern | I: AM-GM equality gives specific values of the three sines}]
For x, y, z $\in$ [0, 2$\pi$]. The number of ordered triplets (x, y, z) satisfying 16($\sin^4$x) + 2($\sin^4$y) + 4($\sin^4$z) $-$ 16(sin x) (sin y) (sin z) + 2 = 0.
\begin{qoptions}
\opt{A}{20}
\opt{B}{18}
\opt{C}{16}
\opt{D}{14}
\end{qoptions}
\end{bankq}
\begin{bankq}[id=SA09-P2-MATH-S3-Q04, source={Score Advanced - Test 09 (29 Apr 2023)}, paper=2, qno=4, type=single, subtopic={General solutions}, answer={6}, answer_src=key, difficulty=5, diff_parts={K1 S2 I2 C1 F0}, diff_note={concepts: solutions  sin x + sin y = sin(x+y) on the diamond abs x + abs y = 1 | steps: identity gives 4 sin(x/2) sin(y/2) sin((x+y)/2) = 0; so x = 0 or y = 0 or x + y = 0 mod 2 pi; intersect with the diamond; count points | I: converting the sum minus sine into a product of sines}]
The number of pairs ordered (x, y) satisfying the equation sin x + sin y = sin (x + y) and $|$x$|$ + $|$y$|$ = 1 is equal to
\end{bankq}
\begin{bankq}[id=SA10-P1-MATH-S3-Q06, source={Score Advanced - Test 10 (06 May 2023)}, paper=1, qno=6, type=numerical, subtopic={General solutions}, answer={4.5}, answer_src=key, difficulty=6, diff_parts={K1 S2 I2 C1 F1}, diff_note={concepts: sum of solutions of a trigonometric equation | steps: reduce to a standard form; solve in the interval; sum solutions}]
If the sum of all the solutions of the equation\par
$\left(2\sec^2\text{x} - \cos\dfrac{16}{9}\text{x}\right)\sec\text{x}$\par $- \left(\cos\dfrac{\text{x}}{3} + 2\tan\text{x}.\sec\text{x}\right)\tan\text{x} = \sin\dfrac{\text{x}}{3}$; x $\in$ [$-$9$\pi$, 9$\pi$] is k$\pi$, then k is equal to :
\end{bankq}
\begin{bankq}[id=SA10-P2-MATH-S3-Q07, source={Score Advanced - Test 10 (06 May 2023)}, paper=2, qno=7, type=numerical, subtopic={Graphs}, answer={4}, answer_src=key, difficulty=6, diff_parts={K1 S2 I2 C1 F1}, diff_note={concepts: number of solutions with inverse tangent and root | steps: range of each side; compare; count roots | I: bounded ranges}]
The number of solution(s) of the equation $\left|\tan^{-1}2\,|\text{x}|\right|$ = $\sqrt{16\text{x}^4 - 8\text{x}^2 + 1}$, is(are)
\end{bankq}
\begin{bankq}[id=SA11-P1-MATH-S2-Q09, source={Score Advanced - Test 11 (09 May 2023)}, paper=1, qno=9, type=single, subtopic={Graphs}, answer={C}, answer_src=key, difficulty=6, diff_parts={K1 S3 I2 C1 F0}, diff_note={concepts: min of two functions on an interval | steps: compare sec^2 and the other curve; find the crossing; range; extrema; match | I: pointwise minimum of two curves}]
Let $f$(x) = min\{sec$^2\pi$x, 3$-\sqrt{2}$ sin$\pi$x\}, x $\in$ $\left(0, \dfrac{1}{2}\right)$\par
{\footnotesize\setlength{\tabcolsep}{1.5pt}\noindent\begin{tabular}{|c|L{4.3cm}|c|c|}
\hline
 & \multicolumn{1}{c|}{List-I} & & \multicolumn{1}{c|}{List-II}\\ \hline
(A) & Number of integers in the range of $f$(x) is & (P) & 0\\ \hline
(B) & Number of points of local extremum of $f$(x) is & (Q) & 2\\ \hline
(C) & Number of values of x where $f$(x) is an integer is & (R) & 1\\ \hline
(D) & Number of solutions of the equation $f$(x) = 1 is & (S) & 3\\ \hline
 & & (T) & 4\\ \hline
\end{tabular}}\par
\begin{qoptions}
\opt{A}{A $\longrightarrow$ S;B $\longrightarrow$ R;C $\longrightarrow$ Q;D $\longrightarrow$ P}
\opt{B}{A $\longrightarrow$ R;B $\longrightarrow$ Q;C $\longrightarrow$ Q;D $\longrightarrow$ S}
\opt{C}{A $\longrightarrow$ R;B $\longrightarrow$ R;C $\longrightarrow$ R;D $\longrightarrow$ R}
\opt{D}{A $\longrightarrow$ P;B $\longrightarrow$ S;C $\longrightarrow$ Q;D $\longrightarrow$ R}
\end{qoptions}
\end{bankq}
\begin{bankq}[id=SA11-P1-MATH-S3-Q07, source={Score Advanced - Test 11 (09 May 2023)}, paper=1, qno=7, type=numerical, subtopic={Trig inequalities}, answer={8}, answer_src=key, difficulty=6, diff_parts={K1 S2 I2 C1 F1}, diff_note={concepts: trigonometric equation with a maximum | steps: rewrite as a sum with bound; equality case; values of x y; compute | I: sum of two terms at most 3}, also_in={Score Advanced - Test 29 (01 May 2024) P1 Q7}]
Let x, y $\in$ R such that cos x cos y + 2sin y + 2sin x cos y = 3 then the value of tan$^2$ x + 5tan$^2$ y is
\end{bankq}
\begin{bankq}[id=SA13-P1-MATH-S2-Q07, source={Score Advanced - Test 13 (13 May 2023)}, paper=1, qno=7, type=numerical, subtopic={General solutions}, answer={3}, answer_src=key, difficulty=8, diff_parts={K1 S4 I2 C1 F1}, diff_note={concepts: equation of sum of cubes type | steps: a^3+b^3+c^3 = 3abc with a=20 sin x, b=15 cos x, c = 12 ; factor; solve; count; identify the cube terms; sum of cubes minus 3abc; case a+b+c=0 or a=b=c; solve the trigonometric system; count solutions in 0 to 2pi; check periodicity | I: sum of cubes minus three times product}]
Number of solution(s) of the equation 8000 sin$^3$x + 3375 cos$^3$ x + 1728 = 5400 sin 2x in (0, 2$\pi$) is(are)\par
\end{bankq}
\begin{bankq}[id=SA16-P1-MATH-S1-Q04, source={Score Advanced - Test 16 (20 May 2023)}, paper=1, qno=4, type=single, subtopic={General solutions}, answer={A}, answer_src=key, difficulty=6, diff_parts={K1 S3 I1 C2 F0}, diff_note={concepts: number  solutions of a product of trigonometric expressions equal 40 | steps: first factor 2 sin^2 x - 2 sqrt3 sin x cos x + 1 = 2 - cos 2x range; second factor 9 - 4 cos^2 x - 4 cos x range; product maximum 40 only at specific values; solve; count in [0,20] | C: bounding both factors and counting points where each reaches an extreme}]
If\par\resizebox{\linewidth}{!}{$\left(2\sin^2x-2\sqrt3\sin x\cos x+1\right)\left(9-4\cos^2x-4\cos x\right)=40$}\par then number of values of $x$ in $[0,20]$ is
\begin{qoptions}
\opt{A}{3}
\opt{B}{4}
\opt{C}{5}
\opt{D}{6}
\end{qoptions}
\end{bankq}
\begin{bankq}[id=SA17-P2-MATH-S3-Q03, source={Score Advanced - Test 17 (25 May 2023)}, paper=2, qno=3, type=numerical, subtopic={General solutions}, answer={25.00}, answer_src=key, difficulty=7, diff_parts={K1 S3 I1 C2 F1}, diff_note={concepts: sum  solutions of a trigonometric equation in 0 to 2 pi | steps: rewrite left side as sin x + cos x times a quadratic in tan x; right side sqrt(9 + 9 sin 2x) = 3 abs(sin x + cos x); divide; quadratic in tan x; solve with roots; sum all solutions in the interval | C: reducing a product equation with absolute value to quadratic cases}, also_in={Score Advanced - Test 35 (15 May 2024) P2 Q3}]
If sum of the values of $x$ satisfying the equation\par
$\sqrt2\sin\left(x+\frac\pi4\right)\left(\sec^2x-(1+\sqrt3)\tan x+2+\sqrt3\right)=\sqrt{9+18\sin x\cos x}$\par
in $[0,2\pi]$ is $\left(\frac pq\right)\pi$, where $p$ and $q$ are co-prime, then $|p-q|$ is equal to
\end{bankq}
\begin{bankq}[id=SA19-P1-MATH-S3-Q07, source={Score Advanced - Test 19 (04 Apr 2024)}, paper=1, qno=7, type=numerical, subtopic={General solutions}, answer={1}, answer_src=key, difficulty=5, diff_parts={K1 S2 I1 C1 F1}, diff_note={concepts: number  solutions of a logarithmic trigonometric equation on 0 to pi | steps: simplify the nested logarithms to constants; equation in log base sqrt3 of tan x; solve for tan x values; count x in [0,pi] with domain tan x > 0}]
Number of solutions of $x\in[0,\pi]$ satisfying the equation\par
$\left(\log_{\sqrt3}\tan x\right)\left(\sqrt{\dfrac{\log_{\sqrt3}\left(\sqrt3\right)^3}{\log_{\sqrt3}\sqrt3}+\dfrac{\log_{\sqrt3}3}{\log_{\sqrt3}\tan x}}\right)=-1$ is
\end{bankq}
\begin{bankq}[id=SA23-P1-MATH-S2-Q10, source={Score Advanced - Test 23 (14 Apr 2024)}, paper=1, qno=10, type=multi, subtopic={Graphs}, answer={A,B}, answer_src=key, difficulty=8, diff_parts={K1 S3 I2 C2 F1}, diff_note={concepts: values of a for exactly one solution of 1 - cos x = (sqrt3/2) abs x + a | steps: define h(x) = 1 - cos x - (sqrt3/2) x on (0,pi); h(0) = 0; h' = sin x - sqrt3/2 zero at pi/3 and 2pi/3; local min and max values; horizontal line a meets once; ranges of a | I: sign pattern of the derivative gives min and max | C: evaluating h at critical points and endpoints}]
Possible value(s) of ``a'' if equation $1-\cos x=\dfrac{\sqrt3}{2}|x|+a$, x $\in$ (0, $\pi$), has exactly one solution is/are
\begin{qoptions}
\opt{A}{$\dfrac12-\dfrac{\pi}{2\sqrt3}$}
\opt{B}{$\dfrac32-\dfrac{\pi}{\sqrt3}$}
\opt{C}{$1-\dfrac{\sqrt3\pi}{4}$}
\opt{D}{$1-\dfrac1{\sqrt2}-\dfrac{\sqrt3\pi}{8}$}
\end{qoptions}
\end{bankq}
\begin{bankq}[id=SA25-P1-MATH-S1-Q03, source={Score Advanced - Test 25 (21 Apr 2024)}, paper=1, qno=3, type=single, subtopic={General solutions}, answer={C}, answer_src=key, difficulty=5, diff_parts={K1 S3 I1 C1 F0}, diff_note={concepts: number  solutions of sqrt(tan) = 2 sin | steps: square: tan theta = 4 sin^2; sin cos relation; tan >= 0 in first and third quadrants; solve sin theta (1 - 4 sin theta cos theta ); count in [0,2pi]}, also_in={Score Advanced - Test 26 (21 Apr 2024) P1 Q3}]
Number of solutions of $\sqrt{\tan\theta}=2\sin\theta$, $\theta\in[0,2\pi]$ is equal to
\begin{qoptions}
\opt{A}{2}
\opt{B}{4}
\opt{C}{5}
\opt{D}{7}
\end{qoptions}
\end{bankq}
\begin{bankq}[id=SA27-P1-MATH-S2-Q12, source={Score Advanced - Test 27 (24 Apr 2024)}, paper=1, qno=12, type=multi, subtopic={Trig inequalities}, answer={A,B,C}, answer_src=key, difficulty=5, diff_parts={K1 S2 I1 C1 F1}, diff_note={concepts: conditions on a  b for exactly one solution of a quadratic in sin x | steps: let s = sin x with s in (0,1]; the quadratic s^2 - a s + b = 0 must have exactly one root in (0,1] (each s in (0,1) gives two x values in (0,pi) except s = 1); analyse cases of roots; derive conditions}]
If $\sin^2$ x $-$ a$\sin$ x + b = 0 has only one solution in (0, $\pi$) then which of the following statements are correct?
\begin{qoptions}
\opt{A}{a $\in$ ($-\infty$, 1] $\cup$ [2, $\infty$)}
\opt{B}{b $\in$ ($-\infty$, 0] $\cup$ [1, $\infty$)}
\opt{C}{a = 1 + b}
\opt{D}{a = b}
\end{qoptions}
\end{bankq}

\section*{SRGBT papers}
\begin{bankq}[id=SRGBT03-P1-MATH-S1-Q42, source={SRGBT - Test 03 (08 Aug 2021)}, paper=1, qno=42, type=multi, subtopic={General solutions}, answer={A,C}, answer_src=key, difficulty=4, diff_parts={K1 S2 I1 C0 F1}, diff_note={concepts: values of 4 cos t + 3 sin t given 2 cos t + sin t = 1 | steps: solve the first equation by half angle; two solutions; evaluate the second expression}]
If 2 cos $\theta$ + sin $\theta$ = 1, then the value of 4 cos $\theta$ + 3 sin $\theta$ can be equal to :
\begin{qoptions}
\opt{A}{3}
\opt{B}{$-$ 5}
\opt{C}{$\dfrac75$}
\opt{D}{$-$ 4}
\end{qoptions}
\end{bankq}
\begin{bankq}[id=SRGBT03-P1-MATH-S1-Q49, source={SRGBT - Test 03 (08 Aug 2021)}, paper=1, qno=49, type=multi, subtopic={General solutions}, answer={A,B}, answer_src=key, difficulty=6, diff_parts={K1 S3 I1 C1 F1}, diff_note={concepts: sum of solutions of cos 4x + 6 = 7 cos 2x | steps: quadratic in cos 2x; cos 2x = 1; x = n pi; n up to 99; sum 4950 pi}]
Sum of all the solutions of the equation cos 4x + 6 = 7 cos 2x in [0, 314] is k$\pi$. then :
\begin{qoptions}
\opt{A}{sum of digits of k is 18}
\opt{B}{k is divisible by 25}
\opt{C}{sum of digits of k is 15}
\opt{D}{k is divisible by 7}
\end{qoptions}
\end{bankq}
\begin{bankq}[id=SRGBT03-P1-MATH-S1-Q50, source={SRGBT - Test 03 (08 Aug 2021)}, paper=1, qno=50, type=multi, subtopic={General solutions}, answer={A,B}, answer_src=key, difficulty=5, diff_parts={K1 S2 I1 C1 F1}, diff_note={concepts: solutions of sin x + sin y = sin(x+y) on the square abs(x) + abs(y) = 1 | steps: sum to product gives three cases; x = 0 or y = 0 points; x + y = 0 points; count N}]
If N represents the number of real solutions of the equation sin x + sin y = sin(x+y) and $|$x$|$ + $|$y$|$ =1 simultaneously, then N is greater than
\begin{qoptions}
\opt{A}{2}
\opt{B}{4}
\opt{C}{6}
\opt{D}{8}
\end{qoptions}
\end{bankq}
\begin{bankq}[id=SRGBT03-P1-MATH-S1-Q58, source={SRGBT - Test 03 (08 Aug 2021)}, paper=1, qno=58, type=multi, subtopic={General solutions}, answer={3}, answer_src=key, difficulty=6, diff_parts={K1 S3 I1 C1 F1}, diff_note={concepts: solutions of sec t + cosec t = sqrt15 | steps: set u = sin t + cos t; quadratic in u; range of u on (0, pi); sign of sin cos; count roots}]
Number of real roots of the equation sec$\theta$ + cosec$\theta$ = $\sqrt{15}$ lying between 0 and $\pi$ is :
\end{bankq}
\begin{bankq}[id=SRGBT03-P1-MATH-S1-Q59, source={SRGBT - Test 03 (08 Aug 2021)}, paper=1, qno=59, type=multi, subtopic={General solutions}, answer={2}, answer_src=key, difficulty=4, diff_parts={K1 S2 I1 C0 F1}, diff_note={concepts: product of three linear factors in sin t equal to 6 | steps: product increasing in sin t; equals 6 at sin t = -1; count solutions in the interval}]
The number of real solutions of the equation ($\sin\theta$+2)($\sin\theta$+3)($\sin\theta$+4)=6 in $\left[-\pi,\dfrac{5\pi}2\right)$ is :
\end{bankq}
\begin{bankq}[id=SRGBT10-P1-MATH-S1-Q44, source={SRGBT - Test 10 (26 Sep 2021)}, paper=1, qno=44, type=multi, subtopic={General solutions}, answer={A,B}, answer_src=key, difficulty=6, diff_parts={K2 S2 I2 C0 F1}, diff_note={concepts: tan x tan y + tan y tan z + tan z tan x = 1 when x + y + z = pi/2, equality in a^2 + b^2 + c^2 >= ab + bc + ca | steps: set z = pi/2 - (x + y) so cot(x + y) = tan z; sum of squares = 1 = sum of products; equality forces x = y = pi/6 | I: rewriting cot(x + y) as tan of the third angle}]
If $\tan^2x+\tan^2y+\cot^2(x+y)=1$, $x,y\in\left(0,\dfrac\pi4\right)$, then :
\begin{qoptions}
\opt{A}{x can only be $\dfrac\pi6$}
\opt{B}{y can only be $\dfrac\pi6$}
\opt{C}{x can take more than 1 value}
\opt{D}{none of these}
\end{qoptions}
\end{bankq}
\begin{bankq}[id=SRGBT10-P1-MATH-S1-Q48, source={SRGBT - Test 10 (26 Sep 2021)}, paper=1, qno=48, type=multi, subtopic={General solutions}, answer={A,B,C}, answer_src=key, difficulty=4, diff_parts={K1 S2 I0 C1 F1}, diff_note={concepts: radical trigonometric equation with domain conditions | steps: write 1 + sin 2x = (sin x + cos x)^2; square carefully keeping cos 2x >= 0; solutions sin x + cos x = 0 and x = 2n pi}]
$\sqrt{\cos2x}+\sqrt{1+\sin2x}=2\sqrt{\sin x+\cos x}$ \ if
\begin{qoptions}
\opt{A}{$\sin x+\cos x=0$}
\opt{B}{$x=2n\pi$}
\opt{C}{$x=n\pi-\dfrac\pi4$}
\opt{D}{$x=2n\pi\pm\cos^{-1}\left(\dfrac{-1}5\right)$}
\end{qoptions}
\end{bankq}
\begin{bankq}[id=SRGBT10-P1-MATH-S1-Q50, source={SRGBT - Test 10 (26 Sep 2021)}, paper=1, qno=50, type=multi, subtopic={General solutions}, answer={A,B}, answer_src=key, difficulty=3, diff_parts={K1 S2 I0 C0 F1}, diff_note={concepts: comparing sin and cos 2x on one period | steps: cos 2x = 1 - 2 sin^2 x; solve the inequality in sin x; intervals around n pi of width pi/6 and the isolated points n pi + pi/2}]
The solution set of $|\sin x|\le|\cos2x|$ contains
\begin{qoptions}
\opt{A}{$\displaystyle\bigcup_{n\in\mathrm I}\left\{\left[n\pi-\dfrac\pi6,n\pi+\dfrac\pi6\right]\right\}$}
\opt{B}{$\displaystyle\bigcup_{n\in\mathrm I}\left\{n\pi+\dfrac\pi2\right\}$}
\opt{C}{$\displaystyle\bigcup_{n\in\mathrm I}\left\{\left[n\pi-\dfrac\pi8,n\pi+\dfrac\pi8\right]\right\}$}
\opt{D}{$\displaystyle\bigcup_{n\in\mathrm I}\left\{\left[n\pi-\dfrac\pi4,n\pi+\dfrac\pi4\right]\right\}$}
\end{qoptions}
\end{bankq}
\begin{bankq}[id=SRGBT10-P1-MATH-S1-Q52, source={SRGBT - Test 10 (26 Sep 2021)}, paper=1, qno=52, type=multi, subtopic={General solutions}, answer={13}, answer_src=key, difficulty=4, diff_parts={K1 S2 I0 C1 F1}, diff_note={concepts: sin x cos y = 1 needs both factors at +-1 | steps: sin x = cos y = 1 or both -1; list x and y values; count pairs inside x^2 + y^2 <= 9 pi^2 so m = 3 and 78/(2m) = 13}]
The number of ordered pair (x,y) satisfying equation $\sin x\cos y=1$ and $x^2+y^2\le9\pi^2$ is m the find $\dfrac{78}{2\mathrm m}$
\end{bankq}
\begin{bankq}[id=SRGBT10-P1-MATH-S1-Q53, source={SRGBT - Test 10 (26 Sep 2021)}, paper=1, qno=53, type=multi, subtopic={General solutions}, answer={8}, answer_src=key, difficulty=5, diff_parts={K2 S2 I1 C0 F1}, diff_note={concepts: AM-GM equality case, counting solutions on the intervals | steps: x sin^2 x + 1/x >= 2 sin x for the sign of x; equality needs x sin x = 1 and sin^2 y = 1; count roots of x sin x = 1 and the y values to get 8}]
Total number of ordered pairs (x, y) satisfying $x\left(\sin^2x+\dfrac1{x^2}\right)=2\sin x.\sin^2y$, where $x\in(-\pi,0)\cup(0,\pi)$ and $y\in[0,2\pi]$ is / are
\end{bankq}
\begin{bankq}[id=SRGBT10-P1-MATH-S1-Q57, source={SRGBT - Test 10 (26 Sep 2021)}, paper=1, qno=57, type=multi, subtopic={General solutions}, answer={0}, answer_src=key, difficulty=3, diff_parts={K1 S1 I1 C0 F1}, diff_note={concepts: comparing ranges of both sides | steps: left side is at most sqrt2 while tan^2 + cot^2 >= 2 so no roots}]
Number of roots of the equation $|\sin\theta|+|\cos\theta|=\tan^2\theta+\cot^2\theta$
\end{bankq}
\begin{bankq}[id=SRGBT10-P1-MATH-S1-Q60, source={SRGBT - Test 10 (26 Sep 2021)}, paper=1, qno=60, type=multi, subtopic={General solutions}, answer={2}, answer_src=key, difficulty=3, diff_parts={K1 S2 I0 C0 F1}, diff_note={concepts: solving with sin and cos and excluding cos x = 0 | steps: 1 + sin x = 2 cos^2 x = 2(1 - sin^2 x); sin x = 1/2 (sin x = -1 excluded); x = pi/6 and 5 pi/6 so 2}]
The number of solutions of the equation $\tan x+\sec x=2\cos x$ lying in the interval $[0,2\pi]$ is :
\end{bankq}

\section*{Enthusiast Pre-Score papers}
\begin{bankq}[id=EPS02-P1-MATH-S1-Q03, source={Enthusiast Pre-Score - Test 02 (11 Nov 2020)}, paper=1, qno=3, type=single, subtopic={General solutions}, answer={A}, answer_src=key, difficulty=4, diff_parts={K1 S2 I1 C1 F0}, diff_note={concepts: trigonometric equation with a bound | steps: use AM-GM on sin^2 and cosec^2 terms; equality at specific angles; solve}]
The solution of $4\sin^2x+\tan^2x+\operatorname{cosec}^2x+\cot^2x-6=0$, is (where n $\in$ I) -
\begin{qoptions}
\opt{A}{$n\pi\pm\dfrac{\pi}{4}$}
\opt{B}{$2n\pi\pm\dfrac{\pi}{4}$}
\opt{C}{$n\pi+\dfrac{\pi}{3}$}
\opt{D}{$n\pi-\dfrac{\pi}{6}$}
\end{qoptions}
\end{bankq}
\begin{bankq}[id=EPS02-P1-MATH-S2-Q08, source={Enthusiast Pre-Score - Test 02 (11 Nov 2020)}, paper=1, qno=8, type=multi, subtopic={General solutions}, answer={B,D}, answer_src=key, difficulty=6, diff_parts={K1 S2 I2 C1 F1}, diff_note={concepts: equation of two radicals | steps: each radical is a sum of squares; the equation holds only when both sides equal; solve cos theta; count solutions}]
Let $\sqrt{2\cos^2\theta-4\cos\theta+6}+\sqrt{19-6\cos\theta+3\cos^2\theta}=4+\sin^2\theta+2\cos\theta$ where $0\le\theta\le4\pi$ then :
\begin{qoptions}
\opt{A}{Number of solutions of the equation is 2}
\opt{B}{Number of solution the equation is 3}
\opt{C}{Sum of all solutions is $3\pi$}
\opt{D}{Sum of all solutions is $6\pi$}
\end{qoptions}
\end{bankq}
\begin{bankq}[id=EPS02-P1-MATH-S2-Q12, source={Enthusiast Pre-Score - Test 02 (11 Nov 2020)}, paper=1, qno=12, type=multi, subtopic={General solutions}, answer={A,B,C}, answer_src=key, difficulty=4, diff_parts={K1 S2 I0 C1 F1}, diff_note={concepts: trigonometric equation solutions | steps: solve cos tan x = sin tan x; count in the interval; second; third}]
Which of the following statement(s) is(are) correct ?
\begin{qoptions}
\opt{A}{Number of solutions of the equation cos(tanx) = sin(tanx) in x $\in\left(0,\dfrac\pi2\right)$ is infinite.}
\opt{B}{If the equation in x, sinx + cos(k + x) + cos(k $-$ x) = 2 has a real solution, then $|\sin k|\le\dfrac12$.}
\opt{C}{The equation 4sin2x + cosx = 5 has no real solution.}
\opt{D}{Number of solutions of the equation sinx.cos4x.sin5x = $\dfrac{-1}{2}$ in x $\in\left(-\pi,\dfrac{3\pi}2\right]$ is two.}
\end{qoptions}
\end{bankq}
\begin{bankq}[id=EPS02-P1-MATH-S3-Q06, source={Enthusiast Pre-Score - Test 02 (11 Nov 2020)}, paper=1, qno=6, type=numerical, subtopic={General solutions}, answer={6.00}, answer_src=key, difficulty=7, diff_parts={K2 S2 I2 C1 F1}, diff_note={concepts: inequality holds for all x, quadratic in x | steps: discriminant must be nonpositive; sum of squares form; solve for theta phi; count pairs | I: perfect square forced}]
If $-x^2+3(\sin\theta)x-\cos\phi(x-\sin\theta)\le0\ \forall\ x\in R$, then number of ordered pair ($\theta$,$\phi$), $\theta,\phi\in[0,2\pi]$ is
\end{bankq}
\begin{bankq}[id=EPS02-P2-MATH-S1-Q02, source={Enthusiast Pre-Score - Test 02 (11 Nov 2020)}, paper=2, qno=2, type=multi, subtopic={General solutions}, answer={A,C}, answer_src=key, difficulty=5, diff_parts={K1 S1 I2 C1 F1}, diff_note={concepts: trigonometric equation with modulus | steps: range of | sin 4x}]
The number of values of $\alpha$ in [0, $3\pi$] satisfying the equation $|\sin4x|+|\cos4x|=|\cos2\alpha|$ is A and sum of all values of $\alpha$ is B then which of the following option(s) is/are correct-
\begin{qoptions}
\opt{A}{A + 2B = 7 + 21$\pi$}
\opt{B}{A + 2B = 6 + 21$\pi$}
\opt{C}{2AB = 147$\pi$}
\opt{D}{2AB = 21$\pi$}
\end{qoptions}
\end{bankq}
\begin{bankq}[id=EPS02-P2-MATH-S1-Q03, source={Enthusiast Pre-Score - Test 02 (11 Nov 2020)}, paper=2, qno=3, type=multi, subtopic={Trig inequalities}, answer={A,B,C,D}, answer_src=key, difficulty=6, diff_parts={K1 S2 I2 C1 F1}, diff_note={concepts: inequality with trigonometric bounds | steps: bound each side; AM-GM terms; equality case; solve | I: both sides bounded}]
If $16+\sin^{12}x+\cos^{11}x-\sec^2y-9\operatorname{cosec}^2y\ \ge\ \sec^8z-2\sec z+\log_2\left(5-\cos^2x\right)$ and x,y,z $\in$ [0,2$\pi$] then correct option/s is/are-
\begin{qoptions}
\opt{A}{Number of possible values of x are 2}
\opt{B}{Number of possible ordered pairs(x,y) are 8}
\opt{C}{Number of possible ordered pairs (y,z) are 8}
\opt{D}{Number of possible ordered triplets (x,y,z) are 16}
\end{qoptions}
\end{bankq}
\begin{bankq}[id=EPS02-P2-MATH-S2-Q11, source={Enthusiast Pre-Score - Test 02 (11 Nov 2020)}, paper=2, qno=11, type=single, subtopic={General solutions}, answer={B}, answer_src=key, difficulty=4, diff_parts={K1 S2 I1 C1 F0}, diff_note={concepts: matching with set conditions | steps: set X and the conditions; count elements for each; match}]
\textbf{Answer the following by appropriately matching the lists based on the information given in the paragraph.}\par
If set $X=\left\{x:\dfrac{\sqrt3-1}{\sin x}+\dfrac{\sqrt3+1}{\cos x}=4\sqrt2,\right.$\par $\left.x\in(0,\pi)-\left\{\dfrac\pi2\right\}\right\}$\par
$Y=\left\{x:4|\sin x|-|x-2|=0\right\}$\par
$Z=\left\{f(x):f(x)=\sin(\sin x)+\cos(\sin x)\right\}$\par
$W=\left\{f(x):f(x)=4\sin x\cos x-\sin^2x+3\cos^2x\right\}$\par
\smallskip
\noindent{\footnotesize\setlength{\tabcolsep}{2pt}\begin{tabular}{|L{.95cm}|L{3.2cm}|L{.95cm}|L{1.8cm}|}\hline
\multicolumn{2}{|c|}{List-I} & \multicolumn{2}{c|}{List-II}\\\hline
(I) & Number of elements in set X & (P) & 1\\\hline
(II) & Number of elements in set Y & (Q) & 2\\\hline
(III) & Number of integers in set Z & (R) & 3\\\hline
(IV) & Number of integers in set W & (S) & 4\\\hline
 & & (T) & 5\\\hline
 & & (U) & 6\\\hline
\end{tabular}}\par
Which of the following is the only \textbf{INCORRECT} combination ?
\begin{qoptions}
\opt{A}{(II), (U)}
\opt{B}{(III), (P)}
\opt{C}{(I), (R)}
\opt{D}{(IV), (T)}
\end{qoptions}
\end{bankq}
\begin{bankq}[id=EPS02-P2-MATH-S2-Q12, source={Enthusiast Pre-Score - Test 02 (11 Nov 2020)}, paper=2, qno=12, type=single, subtopic={General solutions}, answer={B}, answer_src=key, difficulty=3, diff_parts={K1 S1 I1 C1 F0}, diff_note={concepts: matching with set conditions | steps: count remaining entries; match}]
\textbf{Answer the following by appropriately matching the lists based on the information given in the paragraph.}\par
If set $X=\left\{x:\dfrac{\sqrt3-1}{\sin x}+\dfrac{\sqrt3+1}{\cos x}=4\sqrt2,\right.$\par $\left.x\in(0,\pi)-\left\{\dfrac\pi2\right\}\right\}$\par
$Y=\left\{x:4|\sin x|-|x-2|=0\right\}$\par
$Z=\left\{f(x):f(x)=\sin(\sin x)+\cos(\sin x)\right\}$\par
$W=\left\{f(x):f(x)=4\sin x\cos x-\sin^2x+3\cos^2x\right\}$\par
\smallskip
\noindent{\footnotesize\setlength{\tabcolsep}{2pt}\begin{tabular}{|L{.95cm}|L{3.2cm}|L{.95cm}|L{1.8cm}|}\hline
\multicolumn{2}{|c|}{List-I} & \multicolumn{2}{c|}{List-II}\\\hline
(I) & Number of elements in set X & (P) & 1\\\hline
(II) & Number of elements in set Y & (Q) & 2\\\hline
(III) & Number of integers in set Z & (R) & 3\\\hline
(IV) & Number of integers in set W & (S) & 4\\\hline
 & & (T) & 5\\\hline
 & & (U) & 6\\\hline
\end{tabular}}\par
Which of the following is the only \textbf{CORRECT} combination ?
\begin{qoptions}
\opt{A}{(IV), (U)}
\opt{B}{(IV), (T)}
\opt{C}{(I), (P)}
\opt{D}{(I), (S)}
\end{qoptions}
\end{bankq}

\section*{JEE Advanced (2021 Series) papers}
\begin{bankq}[id=JA01-P1-MATH-S3-Q06, source={JEE Advanced (2021 Series) - Test 01 (24 Jan 2021)}, paper=1, qno=6, type=numerical, subtopic={General solutions}, answer={12.56}, answer_src=key, difficulty=6, diff_parts={K1 S3 I1 C1 F1}, diff_note={concepts: sum x values satisfying a trigonometric equation with tangents | steps: let a = tan^2 x and b = tan^2 y; AM-GM gives left side at least 4; right side at most 4; equality forces a = b = 1 and sin^2(x+y) = 1; list x values in interval; add}]
If $(x,y)$ satisfy the equation $\tan^4x+\tan^4y+2\cot^2x\cot^2y=3+\sin^2(x+y)$ sum of all the possible values of $x$ where $x,y\in(0,2\pi)$
\end{bankq}
\begin{bankq}[id=JA02-P1-MATH-S3-Q03, source={JEE Advanced (2021 Series) - Test 02 (04 Feb 2021)}, paper=1, qno=3, type=numerical, subtopic={Trig inequalities}, answer={1.00}, answer_src=key, difficulty=5, diff_parts={K1 S2 I1 C1 F1}, diff_note={concepts: length the interval where a trigonometric function is nonnegative | steps: use product to sum on f; reduce to single sine expression; solve inequality; difference of endpoints}]
Let $f:[0,2]\to\mathbb R$ be the function defined by
\[f(x)=(3-\sin(2\pi x))\sin\!\left(\pi x-\frac\pi4\right)-\sin\!\left(3\pi x+\frac\pi4\right)\]
If $\alpha,\beta\in[0,2]$ are such that $\{x\in[0,2]:f(x)\ge0\}=[\alpha,\beta]$, then the value of $\beta-\alpha$ is \underline{\hspace{1.5cm}}
\end{bankq}

\section*{TOAS papers}
\begin{bankq}[id=TOAS04-P1-MATH-S1-Q04, source={TOAS - Test 04 (15 Mar 2022)}, paper=1, qno=4, type=single, subtopic={Trig inequalities}, answer={C}, answer_src=key, difficulty=5, diff_parts={K2 S2 I1 C1 F0}, diff_note={concepts: trigonometric inequalities, cubic inequality | steps: write cos x cos 2x in terms of cos x; cubic inequality; combine with cos x > 1/3 | I: factor the cubic in cos x}, also_in={TOPS - Test 02 (28 Aug 2022) P1 Q2}]
The complete interval of values of $x$ in $\left(-\dfrac\pi2,\dfrac\pi2\right)$ satisfying the inequations $\cos x\cdot\cos2x>-\dfrac14$ and $3\cos x>1$, is
\begin{qoptions}
\opt{A}{$\left(0,\dfrac\pi3\right)$}
\opt{B}{$\left(-\dfrac\pi2,\dfrac\pi4\right)$}
\opt{C}{$\left(-\dfrac\pi3,\dfrac\pi3\right)$}
\opt{D}{$\left(-\dfrac\pi3,\dfrac{5\pi}{12}\right)$}
\end{qoptions}
\end{bankq}
\begin{bankq}[id=TOAS05-P1-MATH-S2-Q11, source={TOAS - Test 05 (03 Apr 2022)}, paper=1, qno=11, type=single, subtopic={General solutions}, answer={B}, answer_src=key, difficulty=3, diff_parts={K1 S2 I0 C1 F0}, diff_note={concepts: general solutions of trigonometric equations | steps: solve item I; item II; item III; item IV}]
Match List-I with List-II and select the correct answer using the code given below the list.
{\small\setlength{\tabcolsep}{3pt}\renewcommand{\arraystretch}{1.25}\begin{center}\begin{tabular}{|L{0.8cm}|L{4.6cm}|L{0.65cm}|L{0.4cm}|}
\hline
\multicolumn{2}{|c|}{List-I} & \multicolumn{2}{c|}{List-II}\\
\hline
(I) & Number of solutions of $4(\sin^6\theta+\cos^6\theta)=10-18\sin\theta\cos\theta$ in $[0,3\pi]$, is & (P) & 1\\
\hline
(II) & Number of possible pairs $(x,y)$ satisfying $x^2-6x+10-\sin^2y=0$, where $y\in[0,2\pi]$ and $x\in R$, is & (Q) & 2\\
\hline
(III) & Number of integral values of $x$ satisfying $\sin x<\dfrac12$ in $[0,2\pi]$, is & (R) & 3\\
\hline
(IV) & Number of solutions satisfying $\tan x+\sec x=2\cos x$ in $[0,4\pi]$, is & (S) & 4\\
\hline
\multicolumn{2}{|c|}{} & (T) & 5\\
\hline
\multicolumn{2}{|c|}{} & (U) & 6\\
\hline
\end{tabular}\end{center}}
Which of the following is the only \textbf{CORRECT} combination ?
\begin{qoptions}
\opt{A}{I $\to$ Q}
\opt{B}{I $\to$ R}
\opt{C}{II $\to$ P}
\opt{D}{II $\to$ R}
\end{qoptions}
\end{bankq}
\begin{bankq}[id=TOAS05-P1-MATH-S2-Q12, source={TOAS - Test 05 (03 Apr 2022)}, paper=1, qno=12, type=single, subtopic={General solutions}, answer={D}, answer_src=key, difficulty=3, diff_parts={K1 S2 I0 C1 F0}, diff_note={concepts: general solutions of trigonometric equations | steps: solve item I; item II; item III; item IV}]
Match List-I with List-II and select the correct answer using the code given below the list.
{\small\setlength{\tabcolsep}{3pt}\renewcommand{\arraystretch}{1.25}\begin{center}\begin{tabular}{|L{0.8cm}|L{4.6cm}|L{0.65cm}|L{0.4cm}|}
\hline
\multicolumn{2}{|c|}{List-I} & \multicolumn{2}{c|}{List-II}\\
\hline
(I) & Number of solutions of $4(\sin^6\theta+\cos^6\theta)=10-18\sin\theta\cos\theta$ in $[0,3\pi]$, is & (P) & 1\\
\hline
(II) & Number of possible pairs $(x,y)$ satisfying $x^2-6x+10-\sin^2y=0$, where $y\in[0,2\pi]$ and $x\in R$, is & (Q) & 2\\
\hline
(III) & Number of integral values of $x$ satisfying $\sin x<\dfrac12$ in $[0,2\pi]$, is & (R) & 3\\
\hline
(IV) & Number of solutions satisfying $\tan x+\sec x=2\cos x$ in $[0,4\pi]$, is & (S) & 4\\
\hline
\multicolumn{2}{|c|}{} & (T) & 5\\
\hline
\multicolumn{2}{|c|}{} & (U) & 6\\
\hline
\end{tabular}\end{center}}
Which of the following is the only \textbf{CORRECT} combination ?
\begin{qoptions}
\opt{A}{III $\to$ S}
\opt{B}{III $\to$ U}
\opt{C}{IV $\to$ T}
\opt{D}{IV $\to$ S}
\end{qoptions}
\end{bankq}
\begin{bankq}[id=TOAS05-P1-MATH-S3-Q06, source={TOAS - Test 05 (03 Apr 2022)}, paper=1, qno=6, type=numerical, subtopic={General solutions}, answer={0.00}, answer_src=key, difficulty=4, diff_parts={K1 S1 I2 C0 F1}, diff_note={concepts: trigonometric equation by bounding | steps: rewrite with product-to-sum; bound each term to show LHS < 0 | I: bounding argument}]
The number of solution of equation $2\sin2x+\sin x(1-2\cos2x)-4=0$ in the interval $(0,\pi)$ :
\end{bankq}
\begin{bankq}[id=TOAS05-P2-MATH-S1-Q17, source={TOAS - Test 05 (03 Apr 2022)}, paper=2, qno=17, type=single, subtopic={General solutions}, answer={C}, answer_src=key, difficulty=3, diff_parts={K1 S2 I1 C0 F0}, diff_note={concepts: trigonometric equation | steps: tan 3 theta = tan theta + tan 2 theta relation; solve for theta; exclude undefined values | I: use the identity for A + B + C}, also_in={TMAS - Test 01 (09 Jul 2023) P1 Q7; TNAS - Test 03 (28 May 2023) P1 Q7}]
Number of solutions of the equation $\tan\theta+\tan2\theta+\tan3\theta=\tan\theta\tan2\theta\tan3\theta$ in $(0,\pi)$ is-
\begin{qoptions}
\opt{A}{5}
\opt{B}{3}
\opt{C}{2}
\opt{D}{zero}
\end{qoptions}
\end{bankq}
\begin{bankq}[id=TOAS05-P2-MATH-S1-Q18, source={TOAS - Test 05 (03 Apr 2022)}, paper=2, qno=18, type=single, subtopic={General solutions}, answer={C}, answer_src=key, difficulty=2, diff_parts={K1 S1 I1 C0 F0}, diff_note={concepts: exponential trigonometric equation | steps: reduce to sin | x}, also_in={TMAS - Test 01 (09 Jul 2023) P1 Q8; TNAS - Test 03 (28 May 2023) P1 Q8}]
The number of solutions of $2^{\sin|x|}=4^{|\cos x|}$ in $[-\pi,\pi]$ is -
\begin{qoptions}
\opt{A}{0}
\opt{B}{2}
\opt{C}{4}
\opt{D}{6}
\end{qoptions}
\end{bankq}
\begin{bankq}[id=TOAS05-P2-MATH-S1-Q19, source={TOAS - Test 05 (03 Apr 2022)}, paper=2, qno=19, type=single, subtopic={General solutions}, answer={A}, answer_src=key, difficulty=2, diff_parts={K1 S1 I1 C0 F0}, diff_note={concepts: general solution with a minimum | steps: x\^{}2 - 6x + 10 >= 1 so min = 1/sqrt2; solve sin(x - pi/4) = 1/2 | I: bound the quadratic}, also_in={TMAS - Test 01 (09 Jul 2023) P1 Q9; TNAS - Test 03 (28 May 2023) P1 Q9}]
The general solution of the equation $\sin x-\cos x=\min\left\{\dfrac1{\sqrt2},x^2-6x+10\right\}$ where $n\in I$, is-
\begin{qoptions}
\opt{A}{$\dfrac\pi4+n\pi+(-1)^n\dfrac\pi6$}
\opt{B}{$n\pi+(-1)^n\dfrac\pi6-\dfrac\pi4$}
\opt{C}{$2n\pi$}
\opt{D}{$2n\pi+\dfrac\pi2$}
\end{qoptions}
\end{bankq}
\begin{bankq}[id=TOAS19-P1-MATH-S1-Q06, source={TOAS - Test 19 (05 Feb 2023)}, paper=1, qno=6, type=multi, subtopic={Trig inequalities}, answer={C}, answer_src=key, difficulty=3, diff_parts={K1 S2 I0 C0 F1}, diff_note={concepts: trigonometric inequality in an interval | steps: condition sin 3x <= 1/2; solve in the interval; match the options}, also_in={TNAS - Test 01 (05 Feb 2023) P2 Q6}]
The set of all values of $x$ in $[0,\pi]$ for which the function $f(x)=\sqrt{1-2\sin3x}$ takes real values is-
\begin{qoptions}
\opt{A}{$\left[0,\dfrac\pi6\right]\cup\left[\dfrac{5\pi}6,\pi\right]$}
\opt{B}{$\left[0,\dfrac\pi9\right]\cup\left[\dfrac{2\pi}9,\dfrac{7\pi}9\right]\cup\left[\dfrac{8\pi}9,\pi\right]$}
\opt{C}{$\left[0,\dfrac\pi{18}\right]\cup\left[\dfrac{5\pi}{18},\dfrac{13\pi}{18}\right]\cup\left[\dfrac{17\pi}{18},\pi\right]$}
\opt{D}{$\left[0,\dfrac{13\pi}{18}\right]\cup\left[\dfrac{17\pi}{18},\pi\right]$}
\end{qoptions}
\end{bankq}
\begin{bankq}[id=TOAS19-P1-MATH-S2-Q11, source={TOAS - Test 19 (05 Feb 2023)}, paper=1, qno=11, type=single, subtopic={General solutions}, answer={B}, answer_src=key, difficulty=4, diff_parts={K1 S2 I1 C1 F0}, diff_note={concepts: trigonometric equations in intervals | steps: simplify each equation; count the solutions in each interval; match}]
Answer the following by appropriately matching the lists based on the information given in the paragraph

List-I contains some trigonometric equations and List-II provides number of solutions in given interval.
\begin{center}\footnotesize\setlength{\tabcolsep}{2pt}
\begin{tabular}{|L{0.8cm}L{3.8cm}|L{0.6cm}L{1.9cm}|}\hline
\multicolumn{2}{|c|}{List-I}&\multicolumn{2}{c|}{List-II}\\\hline
(I)&$2\cot2x-3\cot3x=\tan2x$&(P)&0 in $\left(0,\dfrac\pi2\right)$\\\hline
(II)&$\sin^{10}x+\cos^{10}x=\dfrac{29}{16}\cos^42x$&(Q)&2 in $(0,\pi)$\\\hline
(III)&$\cos^42x+2\sin^22x=17(\cos x+\sin x)^8$&(R)&4 in $(0,\pi)$\\\hline
(IV)&$\sin\!\left(x+\dfrac\pi9\right)=\cos\!\left(2x-\dfrac{59\pi}{180}\right)$&(S)&4 in $(0,2\pi)$\\\hline
&&(T)&2 in $\left(-\dfrac\pi2,\dfrac\pi2\right)$\\\hline
&&(U)&2 in $\left(\dfrac\pi2,\pi\right)$\\\hline
\end{tabular}\end{center}
Which of the following is the only correct combination ?
\begin{qoptions}
\opt{A}{I $\to$ T}
\opt{B}{II $\to$ R}
\opt{C}{III $\to$ P}
\opt{D}{IV $\to$ R}
\end{qoptions}
\end{bankq}
\begin{bankq}[id=TOAS19-P1-MATH-S2-Q12, source={TOAS - Test 19 (05 Feb 2023)}, paper=1, qno=12, type=single, subtopic={General solutions}, answer={C}, answer_src=key, difficulty=4, diff_parts={K1 S2 I1 C1 F0}, diff_note={concepts: trigonometric equations in intervals | steps: simplify each equation; count the solutions in each interval; match}]
Answer the following by appropriately matching the lists based on the information given in the paragraph

List-I contains some trigonometric equations and List-II provides number of solutions in given interval.
\begin{center}\footnotesize\setlength{\tabcolsep}{2pt}
\begin{tabular}{|L{0.8cm}L{3.8cm}|L{0.6cm}L{1.9cm}|}\hline
\multicolumn{2}{|c|}{List-I}&\multicolumn{2}{c|}{List-II}\\\hline
(I)&$2\cot2x-3\cot3x=\tan2x$&(P)&0 in $\left(0,\dfrac\pi2\right)$\\\hline
(II)&$\sin^{10}x+\cos^{10}x=\dfrac{29}{16}\cos^42x$&(Q)&2 in $(0,\pi)$\\\hline
(III)&$\cos^42x+2\sin^22x=17(\cos x+\sin x)^8$&(R)&4 in $(0,\pi)$\\\hline
(IV)&$\sin\!\left(x+\dfrac\pi9\right)=\cos\!\left(2x-\dfrac{59\pi}{180}\right)$&(S)&4 in $(0,2\pi)$\\\hline
&&(T)&2 in $\left(-\dfrac\pi2,\dfrac\pi2\right)$\\\hline
&&(U)&2 in $\left(\dfrac\pi2,\pi\right)$\\\hline
\end{tabular}\end{center}
Which of the following is the only correct combination ?
\begin{qoptions}
\opt{A}{I $\to$ R}
\opt{B}{II $\to$ T}
\opt{C}{III $\to$ U}
\opt{D}{IV $\to$ R}
\end{qoptions}
\end{bankq}
\begin{bankq}[id=TOAS19-P1-MATH-S3-Q06, source={TOAS - Test 19 (05 Feb 2023)}, paper=1, qno=6, type=numerical, subtopic={General solutions}, answer={50}, answer_src=solved, working={With $c=\cos^4x$, $s=\sin^3x$: $c^3+(-s)^3+1-3c(-s)=0$ so $c-s+1=0$ or $c=-s=1$. The first needs $\cos^4x+1=\sin^3x$, true only for $\sin x=1$, $\cos x=0$; the second is impossible. $x=\pi/2+2k\pi$ in $[0,100\pi]$: $k=0$ to $49$, 50 solutions}, difficulty=7, diff_parts={K1 S3 I2 C1 F1}, diff_note={concepts: sum of cubes identity in a trigonometric equation | steps: put c and s; factor the cube sum; bounds force sin x = 1; second case; count in the interval}]
The number of solutions of the equation $\cos^{12}x+3\cos^4x\sin^3x=\sin^9x-1$ in $x\in[0,100\pi]$ is
\end{bankq}

\section*{TOPS papers}
\begin{bankq}[id=TOPS01-P1-MATH-S3-Q02, source={TOPS - Test 01 (31 Jul 2022)}, paper=1, qno=2, type=numerical, subtopic={General solutions}, answer={3.14}, answer_src=key, difficulty=6, diff_parts={K2 S2 I1 C1 F1}, diff_note={concepts: cube roots, trigonometric equation | steps: let t be the cube root of sin x; solve the cubic; solutions in the interval; mean}]
The arithmetic mean of solutions of equation $\sqrt[3]{\sin x-1}+\sqrt[3]{\sin x}+\sqrt[3]{\sin x+1}=0$, lying in $[0,2\pi]$, is
\end{bankq}
\begin{bankq}[id=TOPS02-P1-MATH-S3-Q04, source={TOPS - Test 02 (28 Aug 2022)}, paper=1, qno=4, type=numerical, subtopic={General solutions}, answer={8.00}, answer_src=key, difficulty=8, diff_parts={K2 S3 I2 C1 F1}, diff_note={concepts: exponential equation, trigonometric equation | steps: let t equal 27 to the sin squared; solve the quadratic; values of sin squared; solutions of x; count in the interval | equation in an interval}]
Number of solution(s) of the equation $27^{\sin^2x}+27^{\cos^2x}=12$ in the interval $[-\pi,\pi]$ is
\end{bankq}
\begin{bankq}[id=TOPS02-P1-MATH-S3-Q05, source={TOPS - Test 02 (28 Aug 2022)}, paper=1, qno=5, type=numerical, subtopic={Trig inequalities}, answer={6.00}, answer_src=key, difficulty=7, diff_parts={K2 S3 I1 C1 F1}, diff_note={concepts: trigonometric inequality, number line | steps: sign analysis of the product; range of cos x; intervals; read off the four endpoints; evaluate}]
If the solution set of inequality $\left(\cos x+\dfrac{\sqrt3}2\right)\left(\cos x-\dfrac12\right)\le0$ in $[0,2\pi]$ is $\left[\dfrac{\alpha\pi}6,\dfrac{\beta\pi}6\right]\cup\left[\dfrac{\gamma\pi}6,\dfrac{\delta\pi}6\right]$, $\alpha,\beta,\gamma,\delta\in I^+$, then the value of $|\beta-\alpha+\delta-\gamma|$ is
\end{bankq}
\begin{bankq}[id=TOPS02-P2-MATH-S2-Q06, source={TOPS - Test 02 (28 Aug 2022)}, paper=2, qno=6, type=numerical, subtopic={General solutions}, answer={2.25}, answer_src=key, difficulty=9, diff_parts={K2 S3 I2 C2 F1}, diff_note={concepts: logarithm, general solution | steps: general solution of the sine; solve for x; apply the interval; values of x; sum | solution in an interval}]
Sum of all solutions of equation $\sin\left(\pi\log_3\left(\dfrac1x\right)\right)=0$ in $[\pi,3\pi]$ is $k$ then $\dfrac k4$ is
\end{bankq}
\begin{bankq}[id=TOPS03-P1-MATH-S2-Q08, source={TOPS - Test 03 (20 Nov 2022)}, paper=1, qno=8, type=multi, subtopic={General solutions}, answer={A,D}, answer_src=key, difficulty=6, diff_parts={K1 S3 I1 C1 F1}, diff_note={concepts: number of solutions of a trigonometric equation with fifth powers | steps: rewrite as sin^5 x - cos^5 x = (sin x - cos x)/(sin x cos x); cancel sin x - cos x; polynomial identity in sin cos; solve; count solutions in each interval}]
The number of solutions of $\sin^5x-\cos^5x=\dfrac1{\cos x}-\dfrac1{\sin x}$ (where $\sin x\ne\cos x$) is :
\begin{qoptions}
\opt{A}{$0$ ; $x\in(0,\pi)-\left\{\dfrac\pi2\right\}$}
\opt{B}{$1$ ; $x\in(0,\pi)-\left\{\dfrac\pi2\right\}$}
\opt{C}{$2$ ; $x\in(0,2\pi)-\left\{\dfrac\pi2,\dfrac{3\pi}2,\pi\right\}$}
\opt{D}{$0$ ; $x\in(0,2\pi)-\left\{\dfrac\pi2,\dfrac{3\pi}2,\pi\right\}$}
\end{qoptions}
\end{bankq}
\begin{bankq}[id=TOPS03-P1-MATH-S2-Q09, source={TOPS - Test 03 (20 Nov 2022)}, paper=1, qno=9, type=multi, subtopic={Graphs}, answer={A,D}, answer_src=key, difficulty=5, diff_parts={K1 S2 I1 C1 F1}, diff_note={concepts: extreme values of (1 + sin t)/(2 - cos t) | steps: set expression to k; a sin t + b cos t = c form; discriminant condition for solvability gives bounds on k; minimum 0 and maximum 4/3}]
If maximum and minimum value of expression $\dfrac{1+\sin\theta}{2-\cos\theta}$ is $M$ and $m$ respectively then :
\begin{qoptions}
\opt{A}{$m=0$}
\opt{B}{$m=\dfrac12$}
\opt{C}{$M=1$}
\opt{D}{$M=\dfrac43$}
\end{qoptions}
\end{bankq}
\begin{bankq}[id=TOPS04-P1-MATH-S4-Q07, source={TOPS - Test 04 (18 Dec 2022)}, paper=1, qno=7, type=numerical, subtopic={General solutions}, answer={5}, answer_src=key, difficulty=8, diff_parts={K2 S3 I2 C1 F1}, diff_note={concepts: exponential equation, trigonometric equation | steps: exponents sum to three; quadratic in the power; sin x values; solutions in the interval; count}]
Number of value of $x$ in $[0,2\pi]$ for which $2^{(2\sin x-1)(\sin x-1)}+2^{(2\sin x+1)(2-\sin x)}=9$, is
\end{bankq}
\begin{bankq}[id=TOPS04-P2-MATH-S1-Q18, source={TOPS - Test 04 (18 Dec 2022)}, paper=2, qno=18, type=single, subtopic={General solutions}, answer={C}, answer_src=key, difficulty=4, diff_parts={K1 S2 I1 C1 F0}, diff_note={concepts: trigonometric equation | steps: sum to product; factor; solutions in the interval}]
Number of solutions of the equation $\cos5x=-\cos3x$ in $x\in[0,\pi]$ is
\begin{qoptions}
\opt{A}{3}
\opt{B}{6}
\opt{C}{5}
\opt{D}{7}
\end{qoptions}
\end{bankq}
\begin{bankq}[id=TOPS04-P2-MATH-S2-Q07, source={TOPS - Test 04 (18 Dec 2022)}, paper=2, qno=7, type=numerical, subtopic={General solutions}, answer={0.00}, answer_src=key, difficulty=5, diff_parts={K1 S2 I1 C1 F1}, diff_note={concepts: trigonometric sum of squares | steps: each term equals one; compatible values; count}]
Number of solution(s) of $x$ in $[0,2\pi]$ for the equation $\displaystyle\sum_{r=1}^{10}\sin^2(rx)=10$ is
\end{bankq}
\begin{bankq}[id=TOPS05-P2-MATH-S2-Q06, source={TOPS - Test 05 (14 Jan 2023)}, paper=2, qno=6, type=multi, subtopic={General solutions}, answer={A,C}, answer_src=key, difficulty=7, diff_parts={K1 S3 I1 C2 F1}, diff_note={concepts: solutions  a trigonometric equation in a long interval | steps: rearrange to factors; sin x =; solutions of tan x = something and cos x =; count in [0,9pi]; sum; check | C: counting solutions of two families over a long interval}]
Let $A$ be the set containing value(s) of $x$ which satisfy the equation $3\sin^2x-\cos x=\sin x(3\cos x-1)$, where $x\in[0,9\pi]$, then mark the correct statement(s).
\begin{qoptions}
\opt{A}{Number of elements in set $A$ is 17}
\opt{B}{Number of elements in set $A$ is 18}
\opt{C}{Sum of elements in set $A$ is $\dfrac{297\pi}4$}
\opt{D}{Sum of elements in set $A$ is $\dfrac{295\pi}4$}
\end{qoptions}
\end{bankq}

\section*{TMAS papers}
\begin{bankq}[id=TMAS01-P1-MATH-S1-Q04, source={TMAS - Test 01 (09 Jul 2023)}, paper=1, qno=4, type=single, subtopic={Graphs}, answer={C}, answer_src=key, difficulty=3, diff_parts={K1 S2 I1 C0 F0}, diff_note={concepts: range of a sine plus cosine | steps: amplitude sqrt(k^2+1)=5; k values; count integers between | I: amplitude form}, also_in={TNAS - Test 03 (28 May 2023) P1 Q4}]
Range of $f(x)=k\sin x+\cos x$ is $[-5,5]$, then total number of integers between the two possible value of `$k$' is -
\begin{qoptions}
\opt{A}{11}
\opt{B}{10}
\opt{C}{9}
\opt{D}{5}
\end{qoptions}
\end{bankq}
\begin{bankq}[id=TMAS01-P1-MATH-S2-Q05, source={TMAS - Test 01 (09 Jul 2023)}, paper=1, qno=5, type=numerical, subtopic={Trig inequalities}, answer={6}, answer_src=key, difficulty=3, diff_parts={K1 S2 I0 C0 F1}, diff_note={concepts: trig inequality | steps: roots of the two factors; sign chart; intervals; sum}, also_in={TNAS - Test 03 (28 May 2023) P1 Q5}]
If the solution set of inequality $\left(\cos x+\dfrac{\sqrt3}2\right)\left(\cos x-\dfrac12\right)\le0$ in $[0,\ 2\pi]$ is $\left[\dfrac{\alpha\pi}6,\dfrac{\beta\pi}6\right]\cup\left[\dfrac{\gamma\pi}6,\dfrac{\delta\pi}6\right]$, $\alpha,\beta,\gamma,\delta\in I^+$, then the value of $|\beta-\alpha+\delta-\gamma|$ is
\end{bankq}
\begin{bankq}[id=TMAS01-P2-MATH-S3-Q04, source={TMAS - Test 01 (09 Jul 2023)}, paper=2, qno=4, type=numerical, subtopic={General solutions}, answer={0}, answer_src=key, difficulty=5, diff_parts={K1 S2 I1 C1 F1}, diff_note={concepts: trig equation bounds | steps: rewrite with sin x; bound using maximum; show no solution; count zero | I: bounding argument}, also_in={TNAS - Test 03 (28 May 2023) P2 Q7}]
The number of solution of equation $2\sin2x+\sin x(1-2\cos2x)-4=0$ in the interval $(0,\pi)$ :
\end{bankq}

\section*{TNAS papers}
\begin{bankq}[id=TNAS01-P2-MATH-S2-Q11, source={TNAS - Test 01 (05 Feb 2023)}, paper=2, qno=11, type=single, subtopic={General solutions}, answer={B}, answer_src=key, difficulty=4, diff_parts={K1 S2 I1 C1 F0}, diff_note={concepts: number of solutions of trigonometric equations | steps: solve each equation in its interval; count; match}]
Answer the following by appropriately matching the lists based on the information given in the paragraph

List-I contains some trigonometric equations and List-II provides number of solutions in given interval.

\begin{center}
{\footnotesize\setlength{\tabcolsep}{2pt}\begin{tabular}{|c|p{4.6cm}|c|p{2.4cm}|}\hline
 & \centering List-I & & \centering List-II\tabularnewline\hline
(I) & $2\cot2x-3\cot3x=\tan2x$ & (P) & $0$ in $\left(0,\frac\pi2\right)$\\\hline
(II) & $\sin^{10}x+\cos^{10}x=\dfrac{29}{16}\cos^42x$ & (Q) & $2$ in $(0,\pi)$\\\hline
(III) & $\cos^42x+2\sin^22x=17(\cos x+\sin x)^8$ & (R) & $4$ in $(0,\pi)$\\\hline
(IV) & $\sin\left(x+\dfrac\pi9\right)=\cos\left(2x-\dfrac{59\pi}{180}\right)$ & (S) & $4$ in $(0,2\pi)$\\\hline
 & & (T) & $2$ in $\left(-\frac\pi2,\frac\pi2\right)$\\\hline
 & & (U) & $2$ in $\left(\frac\pi2,\pi\right)$\\\hline
\end{tabular}}
\end{center}
Which of the following is the only correct combination ?
\begin{qoptions}
\opt{A}{I $\to$ T}
\opt{B}{II $\to$ R}
\opt{C}{III $\to$ P}
\opt{D}{IV $\to$ R}
\end{qoptions}
\end{bankq}
\begin{bankq}[id=TNAS01-P2-MATH-S2-Q12, source={TNAS - Test 01 (05 Feb 2023)}, paper=2, qno=12, type=single, subtopic={General solutions}, answer={C}, answer_src=key, difficulty=4, diff_parts={K1 S2 I1 C1 F0}, diff_note={concepts: number of solutions of trigonometric equations | steps: solve remaining equations; count; match}]
Answer the following by appropriately matching the lists based on the information given in the paragraph

List-I contains some trigonometric equations and List-II provides number of solutions in given interval.

\begin{center}
{\footnotesize\setlength{\tabcolsep}{2pt}\begin{tabular}{|c|p{4.6cm}|c|p{2.4cm}|}\hline
 & \centering List-I & & \centering List-II\tabularnewline\hline
(I) & $2\cot2x-3\cot3x=\tan2x$ & (P) & $0$ in $\left(0,\frac\pi2\right)$\\\hline
(II) & $\sin^{10}x+\cos^{10}x=\dfrac{29}{16}\cos^42x$ & (Q) & $2$ in $(0,\pi)$\\\hline
(III) & $\cos^42x+2\sin^22x=17(\cos x+\sin x)^8$ & (R) & $4$ in $(0,\pi)$\\\hline
(IV) & $\sin\left(x+\dfrac\pi9\right)=\cos\left(2x-\dfrac{59\pi}{180}\right)$ & (S) & $4$ in $(0,2\pi)$\\\hline
 & & (T) & $2$ in $\left(-\frac\pi2,\frac\pi2\right)$\\\hline
 & & (U) & $2$ in $\left(\frac\pi2,\pi\right)$\\\hline
\end{tabular}}
\end{center}
Which of the following is the only correct combination ?
\begin{qoptions}
\opt{A}{I $\to$ R}
\opt{B}{II $\to$ T}
\opt{C}{III $\to$ U}
\opt{D}{IV $\to$ R}
\end{qoptions}
\end{bankq}
\begin{bankq}[id=TNAS01-P2-MATH-S3-Q06, source={TNAS - Test 01 (05 Feb 2023)}, paper=2, qno=6, type=numerical, subtopic={General solutions}, answer={50.00}, answer_src=key, difficulty=6, diff_parts={K1 S2 I2 C1 F1}, diff_note={concepts: trigonometric equation with bounds | steps: rewrite using cos^12 <= 1; bound terms; solve equality case; count in 100 pi | I: extreme values of both sides}]
The number of solutions of the equation $\cos^{12}x+3\cos^4x\sin^3x=\sin^9x-1$ in $x\in[0,100\pi]$ is
\end{bankq}
\begin{bankq}[id=TNAS03-P2-MATH-S2-Q08, source={TNAS - Test 03 (28 May 2023)}, paper=2, qno=8, type=single, subtopic={General solutions}, answer={A}, answer_src=key, difficulty=5, diff_parts={K1 S2 I2 C1 F0}, diff_note={concepts: trigonometric solutions counting | steps: each equation; solutions in the interval; match}]
Match List-I with List-II and select the correct answer using the code given below the list.
\begin{center}
{\footnotesize\setlength{\tabcolsep}{2pt}\begin{tabular}{|c|p{0.5\linewidth}|c|c|}\hline
\multicolumn{2}{|c|}{List-I} & \multicolumn{2}{c|}{List-II}\\\hline
(P) & Number of solutions of $4(\sin^6\theta+\cos^6\theta)=10-18\sin\theta\cos\theta$ in $[0,3\pi]$, is & (1) & $2$\\\hline
(Q) & Number of possible pairs $(x,y)$ satisfying $x^2-6x+10-\sin^2y=0$, where $y\in[0,2\pi]$ and $x\in R$, is & (2) & $3$\\\hline
(R) & Number of integral values of $x$ satisfying $\sin x<\dfrac12$ in $[0,2\pi]$, is & (3) & $4$\\\hline
(S) & Number of solutions satisfying $\tan x+\sec x=2\cos x$ in $[0,4\pi]$, is & (4) & $5$\\\hline
\end{tabular}}
\end{center}
\begin{qoptions}
\opt{A}{$\mathrm P\to2;Q\to1;R\to4;S\to3$}
\opt{B}{$\mathrm P\to1;Q\to2;R\to4;S\to3$}
\opt{C}{$\mathrm P\to1;Q\to2;R\to3;S\to4$}
\opt{D}{$\mathrm P\to2;Q\to1;R\to3;S\to4$}
\end{qoptions}
\end{bankq}
\begin{bankq}[id=TNAS04-P2-MATH-S1-Q04, source={TNAS - Test 04 (18 Jun 2023)}, paper=2, qno=4, type=multi, subtopic={General solutions}, answer={B,C,D}, answer_src=key, difficulty=6, diff_parts={K2 S2 I1 C1 F1}, diff_note={concepts: trigonometric determinant, roots of an equation | steps: expand the determinant; factor; tangent values; count roots on each interval}]
For the equation $\begin{vmatrix}\tan x&\cot x&\cot x\\\cot x&\tan x&\tan x\\\cot x&\cot x&\tan x\end{vmatrix}=0$ which of the following is/are incorrect ?
\begin{qoptions}
\opt{A}{These are four distinct real roots in $[0,2\pi]$}
\opt{B}{$\tan x=\sqrt2$}
\opt{C}{These are six distinct real roots in $[0,2\pi]$}
\opt{D}{Number of distinct real roots in $\left[-\dfrac\pi2,\dfrac\pi2\right]$ is four}
\end{qoptions}
\end{bankq}

\section*{Star Review (Other) papers}
\begin{bankq}[id=STAR01-P1-MATH-S1-Q05, source={Star Review (Other) - Test 01 (12 Sep 2022)}, paper=1, qno=5, type=multi, subtopic={Trig inequalities}, answer={A,B,C,D}, answer_src=key, difficulty=6, diff_parts={K1 S2 I2 C1 F1}, diff_note={concepts: intersection of trigonometric inequality sets | steps: solution of each inequality; intersect two; intersect all three; match intervals}]
If $P=\left\{\theta;\ \sin\theta<\dfrac{-1}{\sqrt2}\right\}$\par
$Q=\left\{\theta;\ \cos\theta<\dfrac{-1}2\right\}$\par
$R=\{\theta;\ \tan\theta>1\}$\par
then , which of the following is/are hold good
\begin{qoptions}
\opt{A}{$P\cap Q\cap R\in\left(\dfrac{5\pi}4,\dfrac{4\pi}3\right)$}
\opt{B}{$P\cap Q\cap R\in\left(\dfrac{-3\pi}4,\dfrac{-2\pi}3\right)$}
\opt{C}{$P\cap Q\in\left(\dfrac{5\pi}4,\dfrac{4\pi}3\right)$}
\opt{D}{$Q\cap R\in\left(\dfrac{-3\pi}4,\dfrac{-2\pi}3\right)$}
\end{qoptions}
\end{bankq}
\begin{bankpassage}[id=STAR01-P1-MATH-S2-G9-10, source={Star Review (Other) - Test 01 (12 Sep 2022)}, questions={9-10}]
If $P=\{x;\ \cos3x=\cos x\ \&\ x\in[0,2\pi]\}$ \& $Q=\{x;\ \tan4x=\tan x\ \&\ x\in[0,2\pi]\}$ then
\end{bankpassage}
\begin{bankq}[id=STAR01-P1-MATH-S2-Q09, source={Star Review (Other) - Test 01 (12 Sep 2022)}, paper=1, qno=9, type=single, subtopic={General solutions}, answer={C}, answer_src=key, difficulty=6, diff_parts={K1 S2 I2 C2 F0}, diff_note={concepts: common solutions of two trigonometric equations | steps: solve the cosine equation; solve the tangent equation; exclusions; intersect}, group=STAR01-P1-MATH-S2-G9-10]
$n(P\cap Q)$ is \_\_\_\_\_\_\_\_\_
\begin{qoptions}
\opt{A}{$1$}
\opt{B}{$2$}
\opt{C}{$3$}
\opt{D}{$4$}
\end{qoptions}
\end{bankq}
\begin{bankq}[id=STAR01-P1-MATH-S2-Q10, source={Star Review (Other) - Test 01 (12 Sep 2022)}, paper=1, qno=10, type=single, subtopic={General solutions}, answer={B}, answer_src=key, difficulty=6, diff_parts={K1 S2 I2 C2 F0}, diff_note={concepts: union of two solution sets | steps: count each; count the intersection; union}, group=STAR01-P1-MATH-S2-G9-10]
$n(P\cup Q)$ is \_\_\_\_\_\_\_\_\_
\begin{qoptions}
\opt{A}{$7$}
\opt{B}{$9$}
\opt{C}{$11$}
\opt{D}{$12$}
\end{qoptions}
\end{bankq}
\begin{bankpassage}[id=STAR01-P1-MATH-S3-G3-4, source={Star Review (Other) - Test 01 (12 Sep 2022)}, questions={3-4}]
For the equation $\log_{\frac12}|\sin x|=2-\log_{\frac12}|\cos x|$
\end{bankpassage}
\begin{bankq}[id=STAR01-P1-MATH-S3-Q03, source={Star Review (Other) - Test 01 (12 Sep 2022)}, paper=1, qno=3, type=numerical, subtopic={General solutions}, answer={2.00}, answer_src=key, difficulty=7, diff_parts={K1 S2 I2 C2 F1}, diff_note={concepts: logarithmic trigonometric equation | steps: convert the logs; get the product of sine and cosine; solve in the interval; sum}, group=STAR01-P1-MATH-S3-G3-4]
The sum of solution(s) in $[0,\pi]$ is $k\pi$, then the value of $k$ is \_\_\_\_\_\_\_\_\_
\end{bankq}
\begin{bankq}[id=STAR01-P1-MATH-S3-Q04, source={Star Review (Other) - Test 01 (12 Sep 2022)}, paper=1, qno=4, type=numerical, subtopic={General solutions}, answer={4.00}, answer_src=key, difficulty=6, diff_parts={K1 S2 I2 C1 F1}, diff_note={concepts: logarithmic trigonometric equation | steps: convert; solve; solutions in the interval; count}, group=STAR01-P1-MATH-S3-G3-4]
The number of solution(s) in $\left[\dfrac{-\pi}2,\dfrac\pi2\right]$ is
\end{bankq}
\begin{bankq}[id=STAR01-P1-MATH-S4-Q09, source={Star Review (Other) - Test 01 (12 Sep 2022)}, paper=1, qno=9, type=numerical, subtopic={General solutions}, answer={4}, answer_src=key, difficulty=7, diff_parts={K1 S2 I2 C2 F1}, diff_note={concepts: exponential trigonometric equation | steps: substitute for the exponents; relation between the two terms; solve; count in the interval}]
The number of solutions of the equation $3^{\sin2x+2\cos^2x}+3^{1-\sin2x+2\sin^2x}=28$ in $x\in[0,2\pi]$ is equal to
\end{bankq}
\begin{bankq}[id=STAR01-P2-MATH-S1-Q05, source={Star Review (Other) - Test 01 (12 Sep 2022)}, paper=2, qno=5, type=single, subtopic={General solutions}, answer={A}, answer_src=key, difficulty=5, diff_parts={K1 S2 I2 C1 F0}, diff_note={concepts: number of solutions comparing bounded functions | steps: range of the left side; range of the right side; equality case}]
The number of elements in the set $S=\left\{x\in R:2\cos\left(\dfrac{x^2+x}6\right)=4^x+4^{-x}\right\}$ is :
\begin{qoptions}
\opt{A}{1}
\opt{B}{3}
\opt{C}{0}
\opt{D}{infinite}
\end{qoptions}
\end{bankq}
\begin{bankq}[id=STAR01-P2-MATH-S1-Q07, source={Star Review (Other) - Test 01 (12 Sep 2022)}, paper=2, qno=7, type=single, subtopic={General solutions}, answer={A}, answer_src=key, difficulty=4, diff_parts={K1 S2 I1 C1 F0}, diff_note={concepts: solving a quartic trigonometric equation | steps: reduce to sine squared; solve the quadratic; general solution}]
If $4\sin^4x+\cos^4x=1$, then possible values of $x$ is equal to -
\begin{qoptions}
\opt{A}{$n\pi$, $n\in I$}
\opt{B}{$n\pi\pm\sin^{-1}\left(\dfrac25\right)$, $n\in I$}
\opt{C}{$n\pi+\dfrac\pi6$, $n\in I$}
\opt{D}{none of these}
\end{qoptions}
\end{bankq}
\begin{bankq}[id=STAR01-P2-MATH-S1-Q08, source={Star Review (Other) - Test 01 (12 Sep 2022)}, paper=2, qno=8, type=single, subtopic={General solutions}, answer={C}, answer_src=key, difficulty=5, diff_parts={K1 S2 I2 C1 F0}, diff_note={concepts: inequality with exponentials and tangent | steps: bound of the exponential factor; bound of the square root; equality cases; count}]
If $2^{\sec^2x}\sqrt{\tan^2y-\tan y+\dfrac12}\le1$, then number of values of $x$ can be, $x\in[0,\ 2\pi]$ -
\begin{qoptions}
\opt{A}{0}
\opt{B}{1}
\opt{C}{3}
\opt{D}{infinite}
\end{qoptions}
\end{bankq}
\begin{bankq}[id=STAR01-P2-MATH-S1-Q09, source={Star Review (Other) - Test 01 (12 Sep 2022)}, paper=2, qno=9, type=single, subtopic={General solutions}, answer={A}, answer_src=key, difficulty=4, diff_parts={K1 S2 I1 C1 F0}, diff_note={concepts: linear trigonometric equation | steps: reduce to the cosine form; angle; general solution}]
Values of $\theta$ satisfying the equation $\left(\sqrt3-1\right)\cos\theta+\left(\sqrt3+1\right)\sin\theta=\left(\sqrt3-1\right)$ is :-
\begin{qoptions}
\opt{A}{$2n\pi+\dfrac{5\pi}6$}
\opt{B}{$2n\pi-\dfrac{5\pi}6$}
\opt{C}{$n\pi\pm\dfrac{5\pi}{12}$}
\opt{D}{$2n\pi\pm\dfrac{5\pi}{12}$}
\end{qoptions}
\end{bankq}
\begin{bankq}[id=STAR01-P2-MATH-S1-Q17, source={Star Review (Other) - Test 01 (12 Sep 2022)}, paper=2, qno=17, type=single, subtopic={General solutions}, answer={D}, answer_src=key, difficulty=5, diff_parts={K1 S2 I2 C1 F0}, diff_note={concepts: trigonometric equation with a bound | steps: rewrite; maximum of the left side; compare with three}]
For $x\in(0,\pi)$, the equation $\sin x+2\sin2x-\sin3x=3$ has
\begin{qoptions}
\opt{A}{infinitely many solutions}
\opt{B}{three solutions}
\opt{C}{one solution}
\opt{D}{no solution}
\end{qoptions}
\end{bankq}
\begin{bankq}[id=STAR01-P2-MATH-S1-Q19, source={Star Review (Other) - Test 01 (12 Sep 2022)}, paper=2, qno=19, type=single, subtopic={General solutions}, answer={D}, answer_src=key, difficulty=6, diff_parts={K1 S2 I2 C2 F0}, diff_note={concepts: number of solutions of a trigonometric identity equation | steps: express in terms of cosine of double angle; solve; count in the interval}]
The number of distinct solutions of equation $\dfrac54\cos^22x+\cos^4x+\sin^4x+\cos^6x+\sin^6x=2$ in the interval $[0,2\pi]$ is
\begin{qoptions}
\opt{A}{2}
\opt{B}{4}
\opt{C}{6}
\opt{D}{8}
\end{qoptions}
\end{bankq}
\begin{bankq}[id=STAR01-P2-MATH-S2-Q02, source={Star Review (Other) - Test 01 (12 Sep 2022)}, paper=2, qno=2, type=numerical, subtopic={General solutions}, answer={16}, answer_src=key, difficulty=8, diff_parts={K1 S3 I2 C2 F1}, diff_note={concepts: trigonometric equation and sum of roots of quadratics | steps: reduce the equation; values of theta; quadratic for each; sum of roots; total}]
Let $S=\{\theta\in(0,2\pi):7\cos^2\theta-3\sin^2\theta-2\cos^22\theta=2\}$. Then, the sum of roots of all the equations $x^2-2\left(\tan^2\theta+\cot^2\theta\right)x+6\sin^2\theta=0\ \theta\in S$, is\_\_\_\_\_\_\_\_.
\end{bankq}
\begin{bankq}[id=STAR01-P2-MATH-S2-Q04, source={Star Review (Other) - Test 01 (12 Sep 2022)}, paper=2, qno=4, type=numerical, subtopic={General solutions}, answer={5}, answer_src=solved, working={The equation factors as $(1-\tan\theta\tan2\theta)(\tan3\theta-\sqrt5)=0$; the first factor gives no solutions, $\tan3\theta=\sqrt5$ has 5 values in $S$}, difficulty=8, diff_parts={K1 S3 I2 C2 F1}, diff_note={concepts: equation from the tangent addition formula | steps: rearrange to a product; tangent of three theta; solutions in the interval; exclusions; count}]
Let $S=\left[-\pi,\dfrac\pi2\right)-\left\{-\dfrac\pi2,-\dfrac\pi4,-\dfrac{3\pi}4,\dfrac\pi4\right\}$. Then the number of elements in the set $A=\left\{\theta\in S:\tan\theta\left(1+\sqrt5\tan(2\theta)\right)=\sqrt5-\tan(2\theta)\right\}$ is \_\_\_\_\_
\end{bankq}
\begin{bankq}[id=STAR01-P2-MATH-S2-Q05, source={Star Review (Other) - Test 01 (12 Sep 2022)}, paper=2, qno=5, type=numerical, subtopic={General solutions}, answer={3}, answer_src=solved, working={$\sin\theta=-\tfrac12$ is common to both: $\theta=\tfrac{7\pi}6,\tfrac{11\pi}6$, sum $3\pi$}, difficulty=6, diff_parts={K1 S2 I2 C1 F1}, diff_note={concepts: common solutions of two trigonometric equations | steps: solve each; intersect; sum; coefficient}]
If the sum of solutions of the system of equations $2\sin^2\theta-\cos2\theta=0$ and $2\cos^2\theta+3\sin\theta=0$ in the interval $[0,2\pi]$ is $k\pi$, then $k$ is equal to \_\_\_\_\_\_.
\end{bankq}
\begin{bankq}[id=STAR02-P2-MATH-S2-Q09, source={Star Review (Other) - Test 02 (10 Oct 2022)}, paper=2, qno=9, type=multi, subtopic={General solutions}, answer={C,D}, answer_src=key, difficulty=5, diff_parts={K1 S2 I1 C1 F1}, diff_note={concepts: trigonometric equation homogeneous in sine and cosine | steps: divide by cosine squared; solve the quadratic; general solution}]
$\sin^2x+2\sin x\cos x-3\cos^2x=0$ if
\begin{qoptions}
\opt{A}{$\tan x=3$}
\opt{B}{$\tan x=-1$}
\opt{C}{$x=n\pi+\pi/4$, $n\in I$}
\opt{D}{$x=n\pi+\tan^{-1}(-3)$, $n\in I$}
\end{qoptions}
\end{bankq}
\end{document}
